% Options for packages loaded elsewhere
\PassOptionsToPackage{unicode}{hyperref}
\PassOptionsToPackage{hyphens}{url}
\PassOptionsToPackage{dvipsnames,svgnames,x11names}{xcolor}
\documentclass[
  11pt,
  letterpaper,
]{article}
\usepackage{xcolor}
\usepackage[margin=1in]{geometry}
\usepackage{amsmath,amssymb}
\setcounter{secnumdepth}{-\maxdimen} % remove section numbering
\usepackage{iftex}
\ifPDFTeX
  \usepackage[T1]{fontenc}
  \usepackage[utf8]{inputenc}
  \usepackage{textcomp} % provide euro and other symbols
\else % if luatex or xetex
  \usepackage{unicode-math} % this also loads fontspec
  \defaultfontfeatures{Scale=MatchLowercase}
  \defaultfontfeatures[\rmfamily]{Ligatures=TeX,Scale=1}
\fi
\usepackage{lmodern}
\ifPDFTeX\else
  % xetex/luatex font selection
\fi
% Use upquote if available, for straight quotes in verbatim environments
\IfFileExists{upquote.sty}{\usepackage{upquote}}{}
\IfFileExists{microtype.sty}{% use microtype if available
  \usepackage[]{microtype}
  \UseMicrotypeSet[protrusion]{basicmath} % disable protrusion for tt fonts
}{}
\makeatletter
\@ifundefined{KOMAClassName}{% if non-KOMA class
  \IfFileExists{parskip.sty}{%
    \usepackage{parskip}
  }{% else
    \setlength{\parindent}{0pt}
    \setlength{\parskip}{6pt plus 2pt minus 1pt}}
}{% if KOMA class
  \KOMAoptions{parskip=half}}
\makeatother
\setlength{\emergencystretch}{3em} % prevent overfull lines
\providecommand{\tightlist}{%
  \setlength{\itemsep}{0pt}\setlength{\parskip}{0pt}}
\usepackage{microtype}
\usepackage{needspace}
\usepackage{etoolbox}
\pretocmd{\section}{\Needspace{5\baselineskip}}{}{}
\pretocmd{\subsection}{\Needspace{4\baselineskip}}{}{}
\setlength{\emergencystretch}{3em}
\allowdisplaybreaks[1]
% Latin Modern Math supplies the standard small set-difference glyph.
\AtBeginDocument{\let\setminus\smallsetminus}
\usepackage{bookmark}
\IfFileExists{xurl.sty}{\usepackage{xurl}}{} % add URL line breaks if available
\urlstyle{same}
% fallback for those not using the hyperref driver hyperxmp:
\makeatletter
\@ifundefined{xmpquote}{\newcommand{\xmpquote}[1]{#1}}{}
\makeatother
\hypersetup{
  pdftitle={Knowing machines with mutual truthfulness and mechanicalness},
  colorlinks=true,
  linkcolor={Maroon},
  filecolor={Maroon},
  citecolor={Blue},
  urlcolor={blue},
  pdfcreator={LaTeX via pandoc}}

\title{Knowing machines with mutual truthfulness and mechanicalness}
\usepackage{etoolbox}
\makeatletter
\providecommand{\subtitle}[1]{% add subtitle to \maketitle
  \apptocmd{\@title}{\par {\large #1 \par}}{}{}
}
\makeatother
\subtitle{A proof of Alexander\textquotesingle s Question 3-1-3}
\author{}
\date{13 October 2026}

\begin{document}
\maketitle
\begin{abstract}
We construct a uniformly computably enumerable family of knowing
machines for every computably enumerable well-founded hierarchy on the
natural numbers. Each machine knows that every machine is truthful and
mechanical, and knows particular program indices for the knowledge
columns of its subordinates. Knowledge is interpreted by consequence in
a theory of modal arithmetic, with arithmetic interpreted in the
standard natural numbers. The result answers Alexander's Question 3-1-3
affirmatively for unrestricted modal formulas and hierarchies of
arbitrary well-founded height. The construction uses a common auxiliary
theory and a recursion-theoretic fixed point. Its soundness proof
combines induction on the hierarchy with induction on finite ordinal
support. Finite patterns of resemblance supply the labelled proof stages
and the finite guards needed to recover source theoremhood. The same
family also satisfies known mechanicalness with parameters and inherits
entire subordinate theories along finite paths.
\end{abstract}

\section{Introduction}\label{introduction}

A knowing machine is an effectively enumerable theory whose theorems are
true in the model determined by the theories of all the machines. A
modal operator expresses each machine's knowledge: an open formula is
known at an assignment when its numerical closure follows from that
machine's theory. Such a machine can reason about arithmetic, about
another machine's knowledge, and about its own knowledge. The existence
problem is to arrange these theories so that their claims about one
another remain true despite this self-reference.

For a fixed formula with one free input variable, its \textbf{knowledge
column} is the set of numbers whose numerical instances the machine
knows. There are two distinct ways to describe a machine's computational
nature. To know its \textbf{mechanicalness} is to know that an index
exists for each unary knowledge column. To know a \textbf{particular
code} is to know the column equation for a specified numerical index.
The distinction is the position of the existential quantifier relative
to the knowledge operator. This distinction matters when a machine
reasons about itself as well as other machines.

Alexander's Question 3-1-3 {[}A, Definition 1.1.1{]} asks for machines
that all know one another's truthfulness and mechanicalness, with
particular codes known along a prescribed hierarchy. The hierarchy is
given by a recursively enumerable well-founded relation: a lower agent
is subordinate to a higher one. The relation need not be transitive, and
its lower ideals need not be decidable. We prove the existence of a
uniformly recursively enumerable family for every such relation. Section
1 states the three required schemas precisely; Section 2 specifies the
weak modal logic and epistemic arithmetic in which they are interpreted.
We use \textbf{r.e.} as an abbreviation for recursively enumerable,
equivalently computably enumerable.

The construction and its soundness proof have different roles. First, a
common auxiliary theory and Kleene's recursion theorem determine all the
theories and their literal program indices. This syntactic construction
is completed before any theory is assumed sound. Second, temporary
models interpret knowledge as raw theorem acceptance together with truth
of the accepted formula. These truth filters make reflection available
while soundness is still being proved. Once every raw theorem predicate
is shown sound, the filters can be removed and knowledge becomes actual
source consequence.

The temporary proof stages attach ordinal labels to the modal operators
of finitely many agents. A cut at a label uses only generators supported
below that label. A \textbf{finite guard} is a finite chain of labels
related by the order-2 resemblance relation; finite-pattern reflection
provides separated copies of such guards in which proofs can be placed.
For each target formula, one finite room is fixed before its term
variants and numerical instances are considered. Individual proof
certificates can then use additional finite rooms. This quantifier order
permits formulas and proof depths of arbitrary finite complexity without
selecting one infinite chain of labels.

Sections 3--5 construct the source family and establish the auxiliary
theory's soundness under a general outside-predicate hypothesis. Section
6 states and derives the exact ordinal input from the cited
pattern-of-resemblance results. Sections 7--14 define the labelled
stages, prove syntactic transport, and reconstruct source-language
models. Sections 15--16 combine induction on ordinal support with
well-founded induction on the hierarchy to prove soundness of the entire
family. Section 17 derives the required schemas and removes the
auxiliary operator. Section 18 gives further consequences and states the
scope of the finite-room argument. All notation and intermediate
constructions used in this proof are introduced here.

\section{1. The theorem}\label{the-theorem}

Let \(\mathcal L\) be first-order arithmetic with modal operators
\(K_i\), \(i\in\mathbb N\), interpreted in Alexander's base logic. Fix
the usual enumeration \((W_e)_{e\in\mathbb N}\) of r.e. subsets of
\(\mathbb N\), with a standard arithmetic formula \(W(e,x)\)
representing membership.

Let \(R\) be any r.e. well-founded relation on \(\mathbb N\). The
convention is that \(jRi\) means that \(j\) is subordinate to \(i\).

\textbf{Theorem.} There is a uniformly r.e. family
\((T_i)_{i\in\mathbb N}\) of closed, base-logically consequence-closed
\(\mathcal L\)-theories such that the standard-arithmetic structure \[
\mathcal M\models K_i\varphi[s]
\quad\Longleftrightarrow\quad
T_i\vdash_{\mathrm{BL}}\varphi^s
\tag{T}
\] satisfies self-doubting Epistemic Arithmetic and all three schemas
below:

\begin{enumerate}
\def\labelenumi{\arabic{enumi}.}
\item
  For all \(i,j\) and all \(\mathcal L\)-formulas \(\varphi\), \[
  \mathcal M\models
  K_i\,\operatorname{ucl}(K_j\varphi\to\varphi).
  \tag{Truth}
  \]
\item
  For all \(jRi\) and all \(\varphi\) with
  \(\operatorname{FV}(\varphi)=\{x\}\), \[
  \mathcal M\models
  \exists e\,K_i\forall x
     \bigl(K_j\varphi\leftrightarrow W(e,x)\bigr).
  \tag{Code}
  \]
\item
  For all \(i,j\) and all such unary \(\varphi\), \[
  \mathcal M\models
  K_i\exists e\,\forall x
     \bigl(K_j\varphi\leftrightarrow W(e,x)\bigr).
  \tag{Mechanical}
  \]
\end{enumerate}

Here \(\varphi^s\) is the result of replacing each free variable by the
canonical numeral for its value under \(s\), and \(\operatorname{ucl}\)
is universal closure. The bound index variable is always chosen fresh.
No transitivity, decidable downward closure, recursive rank, finite
branching, bounded height, modal-depth bound, or restriction on
arithmetic complexity is assumed.

The distinction between (Code) and (Mechanical) is the placement of the
existential quantifier relative to \(K_i\). The construction also gives
the parameter version of (Mechanical) stated in Section 18.1.

\section{2. Logical and computability
preliminaries}\label{logical-and-computability-preliminaries}

\subsection{2.1. The base logic}\label{the-base-logic}

Alexander's modal structures satisfy free-variable independence,
invariance under alphabetic changes of bound variables, and capture-free
variable-for-variable substitution. They do not in general satisfy
arbitrary term substitution through a modal operator. We use his sound
and complete base consequence relation, written \(\vdash\) below. See
{[}A, §1.2, especially Definition 1.2.3 and Theorem 1.2.6{]}.

For closed premises, the finite-certificate form of completeness is: \[
\Sigma\vdash\theta
\quad\Longleftrightarrow\quad
\text{some }\sigma_1,\ldots,\sigma_m\in\Sigma
\text{ make }
\sigma_1\to\cdots\to\sigma_m\to\theta
\text{ base-logically valid}.
\tag{L1}
\] The empty list is allowed. Validity is r.e. in any effective
countable modal language. One can enumerate its certificates using the
ordinary first-order encoding in which each modal formula is a predicate
on its ordered free-variable tuple, with the alphabetic and
weak-substitution axioms BL1 and BL2.

All raw theorem predicates below reject nonclosed formulas. A
consequence-closed theory of closed sentences is invariant under
alphabetic changes, since such changes give base-logically equivalent
sentences.

Write \(\mathrm{PA}^{\mathcal L}\) for ordinary PA with induction
extended to every formula of the displayed modal language. Its induction
instances are \[
\operatorname{ucl}\left[
\varphi(x\mid0)\to
\left(\forall x\bigl(\varphi\to\varphi(x\mid S(x))\bigr)
 \to\forall x\,\varphi\right)\right].
\tag{L2}
\] The remaining PA axioms are the usual arithmetic axioms. This is
Alexander's Definition 1.2.19.

\subsection{2.2. Assigned validity and numerical
instantiation}\label{assigned-validity-and-numerical-instantiation}

For every base-logical validity \(\eta\), include each closed numerical
instance \(\eta^s\) as an \textbf{Assigned Validity} axiom. This is a
substantive axiom scheme: a numerical instance of a base-logical
validity need not itself be valid in every base-logical structure {[}A,
Example 1.2.13{]}.

All standard-arithmetic structures constructed below have the
numerical-assignment property \[
M\models\eta[s]\quad\Longleftrightarrow\quad M\models\eta^s.
\tag{L3}
\] Consequently they satisfy Assigned Validity.

A useful permitted numerical-instantiation argument is the following.
Choose a fresh variable \(y\). The variable-instantiation implication \[
\forall x\,\eta\ \to\ \eta(x\mid y)
\] is base-logically valid. Numerically close this implication,
assigning \(y\) the desired number and closing any remaining free
parameters. Assigned Validity supplies the resulting implication as an
axiom. For several universal quantifiers, iterate this procedure. This
derives the concrete instances used below without assuming a general
modal term-substitution rule.

\subsection{2.3. Arithmetic enumeration}\label{arithmetic-enumeration}

Fix an effective syntax and a standard acceptable enumeration of
programs. Choose \(W(e,x)\) by a finite-computation representation, so
that in standard arithmetic \[
\mathbb N\models W(\bar e,\bar n)
\quad\Longleftrightarrow\quad n\in W_e,
\] and every true concrete instance \(W(\bar e,\bar n)\) is provable in
PA. The latter follows by exhibiting the finite accepting computation
and verifying its arithmetic trace. Only this positive trace property is
used; PA is not required to decide nonmembership.

An r.e. sentence theory has an effective unary-column compiler: given a
formula \(\varphi\) with sole free variable \(x\), a program enumerates
the numbers \(n\) for which the theory proves \(\varphi(x\mid\bar n)\).
Dovetailing finite consequence certificates gives this program. The
compiler is uniform in a program enumerating the theory and the formula.

These elementary computability facts, together with the
\(s\)-\(m\)-\(n\) theorem and Kleene's recursion theorem, are the
computability inputs to the source construction.

\subsection{2.4. Knowing machines and self-doubting Epistemic
Arithmetic}\label{knowing-machines-and-self-doubting-epistemic-arithmetic}

For each operator \(K_i\), the pre-closure knowledge axioms are
universal closures of knowledge of actual base-logical validities,
distribution of knowledge over implication, positive introspection, and
truthfulness. Self-doubting knowledge adds knowledge of each validity,
distribution, and introspection axiom. Self-doubting Epistemic
Arithmetic further includes every modal PA axiom and knowledge of that
axiom {[}A, Definitions 1.2.9 and 1.2.19{]}. The word
\emph{self-doubting} describes the background axiom set: knowledge of
one's own truthfulness is an additional requirement, rather than a
background axiom. The (Truth) schema in Section 1 supplies that
additional requirement for the family constructed here.

A family of theories is a family of knowing entities when its standard
consequence interpretation satisfies these knowledge axioms; it is a
family of knowing machines when each theory is r.e. {[}A, Definitions
1.2.15 and 1.2.17{]}. Our conclusion includes joint enumerability of the
entire family and truth of every accepted sentence in its actual
interpretation. Throughout the proof an \textbf{owner} is the agent
whose modal operator is being interpreted, and a \textbf{holder} is the
agent whose theory contains a displayed axiom. They may be different
agents.

\section{3. The fixed source theories}\label{the-fixed-source-theories}

We now specify the theories whose soundness must be proved. The
\textbf{original agents} are the given agents indexed by natural
numbers. The additional owner \(*\) is an auxiliary agent used to supply
a common theory. The common theory is fixed independently of the
hierarchy; each original theory receives codes for its columns and for
the columns of its direct subordinates. Every program numeral selected
in this construction remains fixed throughout the semantic argument.

Temporarily enlarge the language by one operator \(K_*\). Let \[
I=\{*\}\sqcup\mathbb N,\qquad \mathcal L_*=\mathcal L\cup\{K_*\}.
\] All formula schemes in this section range over the complete enlarged
language, including formulas that contain \(K_*\).

\subsection{3.1. Background and own-prefix
closure}\label{background-and-own-prefix-closure}

For each holder \(a\in I\), let \(B_a\) consist of the following closed
sentences:

\begin{itemize}
\tightlist
\item
  All axioms of \(\mathrm{PA}^{\mathcal L_*}\).
\item
  All Assigned Validity sentences.
\item
  Own validity \(\operatorname{ucl}(K_a\varphi)\), for every actual
  base-logical validity \(\varphi\).
\item
  Own deduction \[
  \operatorname{ucl}\bigl(K_a(\varphi\to\psi)
       \to(K_a\varphi\to K_a\psi)\bigr).
  \]
\item
  Own positive introspection \[
  \operatorname{ucl}(K_a\varphi\to K_aK_a\varphi).
  \]
\item
  All-agent reflection \[
  \operatorname{ucl}(K_b\varphi\to\varphi),\qquad b\in I.
  \]
\end{itemize}

For a set \(X\) of closed generators, define \[
\operatorname{cl}_a(X)=
\operatorname{Cn}_{\mathrm{BL}}
 \{K_a^n\sigma:\sigma\in X,\ n\in\mathbb N\},
\tag{S1}
\] where \(K_a^0\sigma=\sigma\), and \(\operatorname{Cn}_{\mathrm{BL}}\)
here retains closed consequences.

The following syntactic fact will be used repeatedly.

\textbf{Own sentence necessity.} If \(S=\operatorname{cl}_a(X)\), with
\(B_a\subseteq X\), and \(S\vdash\sigma\), then \(S\vdash K_a\sigma\).

\textbf{Proof.} Take a finite certificate with prefixed-generator
premises \(\eta_1,\ldots,\eta_m\) and valid closed implication
\(V=\eta_1\to\cdots\to\eta_m\to\sigma\). The next prefix puts
\(K_a\eta_j\) among the axioms. Own validity gives \(K_aV\), and
repeated own deduction gives \(K_a\sigma\). This uses no soundness
assumption. \(\square\)

\subsection{3.2. The common auxiliary
theory}\label{the-common-auxiliary-theory}

We call the common theory an auxiliary \textbf{core} because all
original theories will inherit it. This is a theory of modal arithmetic;
the ordinal core introduced in Section 6 is the separate structure used
to label proof stages.

For a formula \(\varphi\), a chosen input variable \(x\), and a fresh
variable \(e\ne x\), put \[
\operatorname{MC}_b(\varphi;x)=
\operatorname{ucl}\left[
\exists e\,\forall x
 \bigl(K_b\varphi\leftrightarrow W(e,x)\bigr)\right].
\tag{S2}
\] Thus the other free parameters of \(\varphi\) are universally
quantified outside \(\exists e\). The formula need not have \(x\) as its
only free variable, or even have \(x\) free.

Define \[
U=
\operatorname{cl}_*\left(
 B_*+
 \{\operatorname{MC}_j(\varphi;x):
                  j\in\mathbb N,\ \varphi,\ x\}
 \right).
\tag{S3}
\] The primitive mechanicalness targets are the original agents \(j\).
In particular, this definition includes no primitive
\(\operatorname{MC}_*(\varphi;x)\) generators. Arguments \(\varphi\) may
still contain \(K_*\).

The theory \(U\) is r.e. and independent of \(R\). Fix once and for all
a total computable compiler \(d\) such that, for every unary
\(\varphi\), \[
n\in W_{d(\varphi)}
\quad\Longleftrightarrow\quad
U\vdash\varphi(x\mid\bar n).
\tag{S4}
\] This compiler is chosen before the hierarchy-dependent recursion
fixed point.

\subsection{3.3. The original theories and their fixed
point}\label{the-original-theories-and-their-fixed-point}

Let \(q\) be a program intended to enumerate pairs \((i,\sigma)\) of
original owners and their closed theorems. Uniformly from
\(q,j,\varphi\), obtain a program \(c(q,j,\varphi)\) enumerating \[
\{n:q\text{ enumerates }(j,\varphi(x\mid\bar n))\}.
\]

For a prospective \(q\), set \[
\begin{split}
T_i(q)=\operatorname{cl}_i\bigl(
 B_i
 &+\{C_{*,\varphi}(d(\varphi)):\operatorname{FV}(\varphi)=\{x\}\}\\
 &+\{C_{j,\varphi}(c(q,j,\varphi)):
                    jRi,\ \operatorname{FV}(\varphi)=\{x\}\}
\bigr),
\end{split}
\tag{S5}
\] where \[
C_{b,\varphi}(p)=
\forall x\bigl(K_b\varphi\leftrightarrow W(\bar p,x)\bigr).
\tag{S6}
\] The enumeration of \(R\) supplies positive edge witnesses. All other
admissions, own prefixes, and finite logical consequences are
effectively enumerated. Thus there is a total computable operation
taking \(q\) to a program enumerating all pairs \((i,\sigma)\) with
\(T_i(q)\vdash\sigma\).

Kleene's recursion theorem gives a \(q\) whose enumerated pairs are
exactly those produced by this operation. Fix such a \(q\), and set
\(T_i=T_i(q)\). Then \[
n\in W_{c(q,j,\varphi)}
\quad\Longleftrightarrow\quad
T_j\vdash\varphi(x\mid\bar n).
\tag{S7}
\] The source family \(S_*=U,\ S_i=T_i\) is uniformly r.e. and
consequence-closed.

This completes the source construction before any semantic argument. In
every auxiliary model below, these same theories, source formulas, and
program numerals remain fixed. Masking a semantic predicate does not
rerun the fixed point or change a displayed \(W\)-index.

For convenient notation write \[
p(*,\varphi)=d(\varphi),\qquad
p(j,\varphi)=c(q,j,\varphi)\quad(j\in\mathbb N).
\tag{S8}
\] For each source code generator, its \(p(w,\varphi)\) enumerates
exactly the corresponding full source-theorem column.

\section{4. Filtered models, shape agreement, and
genericity}\label{filtered-models-shape-agreement-and-genericity}

To prove source soundness without assuming it, we first interpret
knowledge using a truth filter. We then prove that retaining source
theories on a lower ideal preserves their soundness under all sound
changes to the remaining predicates. The term \textbf{genericity} below
names this precise preservation property.

\subsection{4.1. Filtered truth}\label{filtered-truth}

Write \(\operatorname{Th}(S)\) for the predicate on closed sentences
that holds exactly when \(S\) proves the sentence in the base logic.

A raw predicate family \(P=(P_a)_{a\in I}\) assigns an alphabetically
invariant predicate to closed \(\mathcal L_*\)-formulas, and rejects
nonclosed formulas. Define a standard-arithmetic filtered structure
\(F(P)\) recursively by \[
F(P)\models K_a\varphi[s]
\quad\Longleftrightarrow\quad
P_a(\varphi^s)\ \land\ F(P)\models\varphi[s].
\tag{F1}
\] This is recursion on formula syntax, so it requires no fixed-point or
soundness assumption. Formula induction gives the base-logic structure
conditions and the numerical-assignment property (L3).

Say that \(P_a\) is \textbf{raw-sound in \(F(P)\)} if \[
P_a(\sigma)\quad\Longrightarrow\quad F(P)\models\sigma
\] for every closed \(\sigma\). When every \(P_a\) is raw-sound, the
truth conjunct in (F1) is redundant. Formula induction then identifies
\(F(P)\) with the unfiltered interpretation using exactly the predicates
\(P_a\). Our task is to prove this when \(P_a=\operatorname{Th}(S_a)\).

Reflection is always true in the filtered model, before raw soundness
has been proved.

\subsection{4.2. Column admissibility and finite query
restrictions}\label{column-admissibility-and-finite-query-restrictions}

Call \(P\) \textbf{column-admissible} if it is alphabetic as above and,
for every owner \(a\), formula \(\varphi\), variable \(x\), and
numerical parameter assignment \(s\), there is an \(e\) such that \[
n\in W_e
\quad\Longleftrightarrow\quad
P_a\bigl(\varphi^{s[x:=n]}\bigr)
\qquad(n\in\mathbb N).
\tag{F2}
\] The actual source family has this property by uniform r.e.
theoremhood. Temporary families need only satisfy the displayed
per-column condition. A uniform computable selection of all its
witnesses is not part of the definition.

The \textbf{shape} of a formula keeps logical constructors, modal
owners, and arithmetic predicate symbols, and forgets terms and
binder-variable names. Let \[
Q(\theta)=
\{(a,\operatorname{sh}(\varphi)):
       K_a\varphi\text{ occurs anywhere in }\theta\}.
\tag{F3}
\] Occurrences below other modal operators are included. This is a
finite hereditary set: retaining an argument shape includes every modal
query occurring inside that shape. Numerical closure, alphabetic change,
and term variants preserve shape and the query set.

Two raw families that agree on every closed formula of every owner/shape
class in \(Q(\theta)\) give the same truth value to \(\theta\), at every
assignment. This follows by formula induction. The universal quantifier
case needs agreement on all numerical instances of an argument shape,
which is why agreement is not restricted to a finite list of literal
subformulas.

A selection depending only on owner and shape preserves column
admissibility. For a fixed numerical column, owner and shape do not vary
with \(n\); the selected column is therefore one fixed old r.e. column
or the empty column. A piecewise choice between two column-admissible
families has the same property. The selector itself need not be
computable.

\subsection{4.3. Whole-ideal genericity}\label{whole-ideal-genericity}

Let \(R^+\) be the transitive closure of \(R\), and write \[
D_R(J)=\{i:i\in J\text{ or }iR^+j\text{ for some }j\in J\}.
\] For a finite \(R^+\)-antichain \(J\), define \(G(J)\) to be this
statement:

\begin{quote}
For every column-admissible \(P\) such that \(P_*=\operatorname{Th}(U)\)
and \(P_i=\operatorname{Th}(T_i)\) for all \(i\in D_R(J)\), if every
original \(P_i\) outside \(D_R(J)\) is raw-sound in \(F(P)\), then the
core and every predicate in the retained ideal are raw-sound in
\(F(P)\).
\end{quote}

No effectiveness of \(D_R(J)\) is required. The source predicates
retained on it are the fixed full theorem predicates.

We will prove \(G(\varnothing)\) directly and the other cases by
well-founded induction. The stronger temporary-model class in (F2) is
needed when a finite formula is reconstructed while an infinite lower
ideal is retained literally.

\section{5. Direct genericity of the auxiliary
theory}\label{direct-genericity-of-the-auxiliary-theory}

\textbf{Core genericity lemma.} Suppose \(P\) is column-admissible,
\(P_*=\operatorname{Th}(U)\), and every original \(P_j\) is raw-sound in
\(F(P)\). Then \(P_*\) is raw-sound there. Thus \(G(\varnothing)\)
holds.

\textbf{Proof.} First verify the unprefixed core generators.

Arithmetic holds in the standard first-order part. Modal induction
follows by ordinary induction on the set of numbers satisfying its
formula in \(F(P)\); the numerical-assignment property identifies the
literal \(0\) and successor instances {[}A, Lemma 1.2.22 and Proposition
1.2.23{]}. Actual validities and Assigned Validity are true by the
base-logic structure conditions and (L3).

For core own validity, if \(\varphi\) is valid then \(\varphi\) is true
at every assignment, and \(U\) contains its numerical instance
\(\varphi^s\). These give both conjuncts of \(K_*\varphi[s]\).

For core deduction, the raw theory \(U\) is closed under modus ponens.
The truth conjuncts of its antecedents give the corresponding
implication of truths. For core introspection, \(U\vdash\varphi^s\)
implies \(U\vdash K_*\varphi^s\) by own sentence necessity. If the inner
filtered modal formula is true, that same formula supplies the outer
truth conjunct. All reflection generators hold directly by (F1). None of
these verifications assumes core raw soundness.

Consider \(\operatorname{MC}_j(\varphi;x)\), with \(j\) original. Fix
values for the outer parameters. Column admissibility supplies an \(e\)
enumerating \[
\{n:P_j(\varphi^{s[x:=n]})\}.
\] Original-owner raw soundness makes
\(F(P)\models K_j\varphi[s[x:=n]]\) equivalent to this raw test. The
chosen \(e\) therefore witnesses the code existential for every \(n\).
Freshness of \(e\) ensures its value does not alter the argument column.
Taking every parameter assignment verifies the universally closed
generator. Arguments containing \(K_*\) cause no extra assumption: the
stipulated original raw soundness already concerns their complete truth
in this same \(F(P)\).

Finally, if a core generator \(\sigma\) is true, every \(K_*^n\sigma\)
is true by induction on \(n\). Its body is true by the preceding step
and is a literal core axiom, hence passes the raw \(U\)-test. All
prefixed generators are therefore true. Base-logical soundness and
finite certificates make every \(U\)-theorem true.

This proves the lemma without a staged core argument. \(\square\)

\section{6. The effective finite-guard ordinal
input}\label{the-effective-finite-guard-ordinal-input}

Let \((\mathrm{Ord};<,\leq_1,\leq_2)\) be the true \textbf{pure}
pattern-of-resemblance structure of order 2. The relations are the
corresponding elementary-substructure relations in the pure language; no
arithmetic operations are added to this ordinal language.

Write \(\upsilon_1=1^\infty\). The entire pure core is the initial
segment \([0,\upsilon_1)\). We use its effective presentation and the
following exact properties.

\textbf{Ordinal lemma.} There is a countably infinite recursive label
domain \(L\), with a recursive well-order \(<\) and recursive relations
\(\leq_1,\leq_2\), isomorphic to this core, satisfying:

\begin{enumerate}
\def\labelenumi{\arabic{enumi}.}
\item
  Both resemblance relations are reflexive and transitive, and \[
  \leq_2\ \subseteq\ \leq_1\ \subseteq\ \leq.
  \]
\item
  \textbf{Connectedness:} if \(p\leq_1r\) and \(p\leq q\leq r\), then
  \(p\leq_1q\).
\item
  \textbf{Finite-pattern reflection:} if \(\alpha\leq_1\beta\), every
  finite \(F\subset[0,\beta)\) has a map \[
  f:F\longrightarrow[0,\alpha)
  \] fixing \(F\cap[0,\alpha)\) and preserving and reflecting equality,
  \(<\), \(\leq_1\), and \(\leq_2\) on \(F\).
\item
  \textbf{Finite tallness:} for every natural \(k\), there are \[
  h_0<\cdots<h_k
  \] with \(h_i\leq_2h_j\) for every \(i<j\). Such a list, beginning at
  \(h_0=\nu\), is called a guard \(G_k(\nu)\). Thus a rank-\(k\) guard
  has \(k+1\) points.
\item
  \textbf{Finite room:} suppose \(G_{k+1}(\lambda)\) is given, \(X\) is
  a finite set with \(x<\lambda\) and \(x\leq_1\lambda\) for all
  \(x\in X\), and \(N\in\mathbb N\). Then there are \(N\) lists \[
  Y_j=(y_{j,0},\ldots,y_{j,k}),\qquad j<N,
  \] each a rank-\(k\) guard, such that every point is below and
  \(\leq_1\)-below \(\lambda\), every point lies above every member of
  \(X\), and \[
  \max Y_i<\min Y_j\qquad(i<j<N).
  \tag{O1}
  \] The assertions for \(X=\varnothing\) and \(N=0\) have their usual
  vacuous meanings.
\end{enumerate}

\textbf{Proof of the ordinal lemma.} First work with the actual core.
The inclusions, transitivity, and connectedness are the
elementary-substructure and respecting-forest properties of the true
pure relations {[}W21, pp.~5--6{]}.

For finite reflection, use the points of \(F\cap\alpha\) as parameters
in the existential closure of the complete finite atomic diagram of
\(F\). Include the negations of false atomic facts and the inequalities
of distinct points. This is a \(\Sigma_1\) assertion in the pure
language containing both resemblance relations, true below \(\beta\).
Reflection at \(\alpha\leq_1\beta\) gives its witnesses below
\(\alpha\). The parameter entries are fixed, and the full diagram gives
exactly the required finite isomorphism. The case \(\alpha=\beta\) uses
the identity. A finite map may be extended arbitrarily to other labels
when a total syntactic relabeling is wanted.

Finite tallness follows from the isominimal realization of every
realizable finite pattern in the core {[}W18, Corollaries 4.8 and 4.10;
W21, p.~6{]}. For each \(k\), a finite increasing pairwise-\(\leq_2\)
pattern of length \(k+1\) is realizable; for example it occurs on the
corresponding finite segment \(\upsilon_2,\ldots,\upsilon_{k+2}\) of the
global sequence in {[}W21, Theorem 1.8{]}. Its isominimal finite copy
lies in the core. This produces each finite guard separately.

For room, write the given guard as \[
\lambda=h_0<h_1<\cdots<h_{k+1}.
\] Apply {[}W21, Lemma 1.7(2){]} with \(\alpha=h_0\), \(\beta=h_{k+1}\),
and \[
Y=\{h_0,\ldots,h_k\}\subset[\alpha,\beta).
\] Every \(y\in Y\) satisfies \(y<_1\beta\). The lemma supplies a copy
of \(Y\) below \(\alpha\), above the finite set \(X\), preserving its
pattern over \(X\); all copied points are \(<_1\alpha\). This gives the
first required rank-\(k\) copy. Enlarge \(X\) by that entire copy and
apply the same lemma again. Repeating finitely many times gives
complete-copy separation (O1). Every enlarged parameter set remains
finite and consists of points \(<_1\alpha\). This includes \(k=0\); the
set being copied is then a singleton. The last guard point \(h_{k+1}\)
is deliberately excluded from \(Y\).

The effective presentation is supplied by the elementary-recursive
characterization of the core {[}W17, §5.1 and Corollary 5.8{]}, using
the corrected semantic characterization in {[}W21, Theorem 4.2 and
Corollary 4.6{]}. The latter repairs an earlier semantic proof, so the
corrected identification is part of the cited input. Order and both
resemblance relations are recursive on the canonical notation domain. If
that domain is an infinite recursive set of numerical codes, enumerate
its members in ordinary numerical order and pull back the relations
along this computable bijection with \(\mathbb N\). This gives the
required recursive presentation on a domain coded by all natural
numbers. Every property above transfers along the isomorphism.
\(\square\)

Only finite witnesses from this presentation enter stage enumeration.
Their existence can be used classically in the proof, and a needed
finite witness can also be found by search through its decidable finite
conditions.

A useful consequence of connectedness is that any set \(C\) of points
each \(\leq_1\)-below the same top \(\lambda\) is coherent in its
increasing order: \[
u,v\in C,\quad u\leq v
\quad\Longrightarrow\quad u\leq_1v.
\tag{O2}
\] In particular this holds for all points of the room copies together
with \(X\).

Within the core, finite guards of unbounded rank cannot all lie in one
fixed proper initial segment {[}W17, introduction; W21, p.~6{]}. Nothing
in the proof chooses an infinite \(\leq_2\)-chain in that core.

\section{7. Labels, recorded syntax, and recording
maps}\label{labels-recorded-syntax-and-recording-maps}

This section introduces the syntax of a temporary labelled proof stage.
The active owners are the finitely many source agents whose soundness is
currently being established. External owners are finitely many owners
outside the retained lower ideal; their old predicates are already
sound. All remaining owners are passive. Active operators receive
ordinal labels. An external query can remain \textbf{live}, with its old
predicate, or become a \textbf{ghost}, whose knowledge predicate is
empty. The choice depends on the argument's syntactic shape. Recording
adds these labels and choices; erasure removes them and returns to the
source language.

Let \(L\) be the effective well-ordered label set, with relations
\(\leq_1,\leq_2\). We use reflexivity of \(\leq_1\), the implication
\(\alpha\leq_1\beta\Rightarrow\alpha\leq\beta\), connectedness of
\(\leq_1\), finite-pattern reflection along \(\leq_1\), and the stated
finite-room property. Write \(G_r(\nu)\) for a list of \(r+1\) points
beginning at \(\nu\), strictly increasing and pairwise \(\leq_2\). No
infinite chain is assumed.

Fix finite sets \(A,E\subseteq\mathbb N\), disjoint from one another,
with the following extension side conditions:

\begin{itemize}
\tightlist
\item
  The core \(*\) is neither active nor external.
\item
  If \(a\in A\) and \(wRa\), then \(w\notin A\cup E\).
\end{itemize}

In the genericity application \(A\) is the finite frontier \(J\), and
\(E\) lies outside its downward ideal, so these conditions follow from
antichainness.

Recorded modal heads are \(K_d\), \(K_b^\nu\), and \(Z_e\), with
erasures \(K_d,K_b,K_e\), respectively. All arithmetic terms and all
numerical constants are unchanged by erasure \(\epsilon\). A recorded
formula is \textbf{allowed} when every indexed owner belongs to \(A\),
every unindexed owner lies outside \(A\), and every ghost owner belongs
to \(E\). It is \textbf{globally admissible} when it is allowed and, at
every indexed occurrence \(K_b^\nu\eta\), every indexed label anywhere
inside \(\eta\) is strictly below \(\nu\). Inspection descends through
unindexed heads, ghosts, Boolean connectives, and quantifiers.

The finite support \(\operatorname{supp}(\eta)\) is the set of all
indexed labels in the whole formula. In particular, labels below live or
ghost heads remain in its support. The notation
\(\operatorname{supp}(\eta)<\nu\) means that every point of this finite
set is strictly below \(\nu\), and is vacuous when the support is empty.
Numerical closure, partial numerical substitution, and alphabetic
variation preserve support and global admissibility.

The shape \(\operatorname{sh}(\phi)\) retains logical constructors,
modal owners, and arithmetic relation symbols, but forgets terms and
binder-variable names. Its rank \(r(\phi)\) is its constructor count.
The query set is \[
Q(\phi)=\{(b,\operatorname{sh}(\psi)):K_b\psi
\text{ is any modal occurrence in }\phi\}.
\] It includes occurrences below other modal heads. It is finite and
hereditary: if \((b,s)\in Q(\phi)\), every query inside \(s\) lies in
\(Q(\phi)\).

Let \(H\) be any finite hereditary set of owner/shape queries. Let
\(\ell(b,s)\in L\) be a label table. Define \(t=t_{H,\ell}\)
recursively, leaving terms, binders and Boolean structure unchanged, and
setting \[
t(K_b\psi)=
\begin{cases}
K_b^{\ell(b,\operatorname{sh}\psi)}t\psi,&b\in A,\\
K_b t\psi,&b\in E,\ (b,\operatorname{sh}\psi)\in H,\\
Z_b t\psi,&b\in E,\ (b,\operatorname{sh}\psi)\notin H,\\
K_b t\psi,&b\notin A\cup E.
\end{cases}
\] The only relevant entries of \(\ell\) for a finite formula or
certificate form a finite table. It can be extended arbitrarily
elsewhere. No globally admissible extension of the table is needed to
establish validity preservation.

This construction has four immediate invariants:

\begin{enumerate}
\def\labelenumi{\arabic{enumi}.}
\tightlist
\item
  \(\epsilon(t\phi)=\phi\).
\item
  \(t\) preserves free variables and commutes with capture-free
  syntactic substitutions, alphabetic changes, and numerical closure.
  This includes the substitutions displayed in modal induction. It is a
  property of the translation, not a claim that these substitutions
  preserve base-logical validity.
\item
  All formulas of one source shape have the same recorded support,
  admissibility condition, and query-label assignments.
\item
  If an external occurrence is kept live, heredity makes its entire
  recorded argument ghost-free.
\end{enumerate}

The proofs are structural inductions. At a modal occurrence, syntactic
substitution or numerical closure does not change the argument shape, so
it does not change the chosen head. The induction hypothesis identifies
the transformed argument. The other constructors are unchanged.

\subsection{Lemma 1: preservation of actual base-logical
validity}\label{lemma-1-preservation-of-actual-base-logical-validity}

Erasure, arbitrary label renaming, and every recording map
\(t_{H,\ell}\) preserve base-logical validity in their forward
directions.

\textbf{Proof.} Pull an arbitrary target modal structure back along the
chosen translation. Interpret a source modal application \(K_b\psi\) by
the target interpretation of its translated modal application, with
translated argument. Because the translation preserves free variables,
irrelevant-variable invariance is inherited. Because it preserves
alphabetic equivalence and commutes with the permitted variable
substitutions, the alphabetic and variable-substitution conditions on
modal interpretations are inherited. The arithmetic structure is
unchanged. Induction on formulas then identifies truth in the
pulled-back structure with truth of the translation. A source validity
is consequently true after translation in every target structure.

Equivalently, in Alexander's first-order encoding, a pure modal formula
names a predicate on its ordered tuple of free variables. Each
translation replaces that predicate name by the name of its translated
modal formula without changing arity. Images of the alphabetic and
variable-substitution axioms BL1/BL2 are instances of the same axioms. A
finite first-order validity derivation can therefore be translated line
by line. Predicate identification under erasure causes no problem.

This argument needs only source sharing by owner and argument shape. It
does \textbf{not} require the recording map to make every formula in the
entire infinite source language globally admissible. Global
admissibility is checked separately on each formula admitted to the
stage.

The converse is unavailable: validity of \(\epsilon\eta\) says nothing
by itself about validity of \(\eta\). In particular,
\(K_e\psi\leftrightarrow Z_e\psi\) is not admitted on account of its
tautological erasure.

\textbf{Numerical-substitution distinction.} The fact that a translation
commutes with numerical closure does not assert that numerical closure
preserves base-logical validity. Assigned-validity rows below supply
numerical instances explicitly. None of the subsequent syntactic
arguments uses unrestricted term substitution through modal operators.

\section{8. The actual stage generator
set}\label{the-actual-stage-generator-set}

The stage theories are generated by explicit axioms, rather than defined
by their intended semantic behavior. Background rows express the logical
and arithmetic rules. Main rows translate an original code axiom while
retaining its source numeral; arithmetic marker antecedents ensure that
every guard point belongs to the row's literal support. Own prefixes
then close each holder's axioms under labelled knowledge. The later
soundness proof must verify exactly this generator set.

For each active holder \(a\in A\), define a set \(\Gamma_a\) of closed
recorded generators by the following rules. Take the recorded modal
language to consist of precisely the allowed heads. Thus the cut
predicates for active holders suffice to interpret every indexed head,
while arbitrary formulas in this language need not be globally
admissible.

\subsection{Background rows}\label{background-rows}

Admit each of the following whenever the entire displayed formula is
globally admissible:

\begin{enumerate}
\def\labelenumi{\arabic{enumi}.}
\tightlist
\item
  Every modal PA axiom in the recorded language, including induction
  with an arbitrary recorded formula.
\item
  \(\eta^s\), for every actual base-logical validity \(\eta\) and every
  numerical assignment to its finitely many free variables.
\item
  \(\operatorname{ucl}(K_a^\nu\eta)\), for every actual validity
  \(\eta\).
\item
  Same-label deduction: \[
  \operatorname{ucl}\bigl(
  K_a^\nu(\eta\to\zeta)\to
  (K_a^\nu\eta\to K_a^\nu\zeta)\bigr).
  \]
\item
  Introspection, for \(\mu<\nu\): \[
  \operatorname{ucl}\bigl(
  K_a^\mu\eta\to K_a^\nu K_a^\mu\eta\bigr).
  \]
\item
  Truthfulness \(\operatorname{ucl}(O\eta\to\eta)\) for every allowed
  indexed, unindexed, or ghost head \(O\).
\item
  Contextual collapse, for \(\mu\leq_1\nu\): \[
  \operatorname{ucl}\bigl(
  C[K_b^\mu\eta]\leftrightarrow C[K_b^\nu\eta]\bigr),
  \tag{CC}
  \] where \(C[-]\) is any one-hole recorded formula context and both
  fillings are globally admissible.
\end{enumerate}

There is no source-sharing requirement on these background rows. In
particular, the two fillings of a contextual-collapse row may assign
different labels to different occurrences of one original source shape.
This is essential: contextual collapse changes one occurrence at a time.

\subsection{Padded main-code rows}\label{padded-main-code-rows}

Let \(a\in A\), and let \(C_{w,\phi}\) be one of \(a\)'s actual source
code generators. By the side conditions, \(w\notin A\cup E\), so its
outer head remains \(K_w\).

Choose one hereditary \(H\) and one source-shared label table \(\ell\)
for the whole source sentence and its distinguished atom \(K_w\phi\).
Require \(tC_{w,\phi}\) to be globally admissible. Inspect the whole
recorded distinguished atom \(t(K_w\phi)\). At each live external
occurrence \(K_e\chi\) in it, impose the following on every indexed
occurrence \(K_b^\nu\beta\) inside \(\chi\):

\begin{itemize}
\tightlist
\item
  There is a displayed guard \(G_{r(\epsilon\chi)}(\nu)\).
\item
  Every label \(\xi\in\operatorname{supp}(\beta)\) satisfies both
  \(\xi<\nu\) and \(\xi\leq_1\nu\).
\end{itemize}

Inspection also descends under ghosts. A ghost occurrence itself imposes
no external guard, but a live external occurrence below it does. The
code head \(K_w\) is inspected and imposes no external guard because
\(w\notin E\).

Let \(B\) be a finite set containing every point of every displayed
guard, including its base. Choose distinct arithmetic tautologies
\(\tau_1,\ldots,\tau_m\), one for each listed point \(h_1,\ldots,h_m\)
of \(B\), with their shapes absent from the following finite forbidden
pool: all subformula shapes of the source code sentence and all
subshapes of the argument shapes in \(H\). During a finite proof
transformation the pool may be enlarged by the finitely many shapes of
that certificate; choosing tags outside this larger pool also satisfies
the admission rule. One explicit infinite supply is \[
\tau_0=(\bot\to\bot),\qquad \tau_{n+1}=(\bot\to\tau_n).
\] Their ranks \(2n+3\) are distinct, so a finite forbidden pool is
avoided by choosing sufficiently large indices. These indices are tags,
not encodings of labels. When labels move, the tags do not move.

Admit the closed padded sentence \[
K_a^{h_1}\tau_1\to
\bigl(K_a^{h_2}\tau_2\to\cdots
(K_a^{h_m}\tau_m\to tC_{w,\phi})\cdots\bigr).
\tag{M}
\] The markers are holder-local. Their arguments are arithmetic and have
empty label support, so padding preserves global admissibility. The
exact support of (M) is \[
B\cup\operatorname{supp}(tC_{w,\phi}).
\] Thus every guard point is literal formula support. There are no
mechanicalness main rows in these extension stages.

The source numeral \(\overline{p(w,\phi)}\) in (M) is unchanged. It is
not replaced by an index computed from \(t\phi\), and it does not encode
a fresh or renamed outside owner. The marker bundle may be empty when no
guards are needed; in that case the padded row is just its main
sentence.

\subsection{Own-prefix closure}\label{own-prefix-closure}

If \(\sigma\in\Gamma_a\), \(a\in A\), and
\(\operatorname{supp}(\sigma)<\nu\), admit \(K_a^\nu\sigma\). Iterate
this rule finitely. Every resulting generator is closed and globally
admissible.

Set \[
\Gamma_a^{<\nu}=
\{\sigma\in\Gamma_a:\operatorname{supp}(\sigma)<\nu\}.
\] The cut predicate on a closed formula \(\eta\) is
\(\Gamma_a^{<\nu}\vdash\eta\), where consequence is Alexander's actual
base-logical consequence. It is false on nonclosed formulas. No
admissibility condition is placed on arbitrary logical conclusions. In
particular, this cut is not defined by intersecting source theoremhood
with an erasure condition.

\subsection{Lemma 2: effectivity of the actual
cut}\label{lemma-2-effectivity-of-the-actual-cut}

The predicates \(\sigma\in\Gamma_a\), \(\sigma\in\Gamma_a^{<\nu}\), and
\(\Gamma_a^{<\nu}\vdash\eta\) are uniformly r.e. in their finite
parameters.

\textbf{Proof.} Enumerate finite syntax, finite masks, finite tables,
finite guard lists, finite marker bundles, and finite prefix trees.
Allowed-head checks, heredity, shape equality, freshness and global
support calculations are finite syntactic checks. Every order,
\(\leq_1\), and \(\leq_2\) assertion has a positive enumerating
certificate. Source-code admissions use the actual source schema and
positive \(R\)-edge certificates. Actual validity has an r.e. proof
system. A cut consequence is witnessed by a finite list of admitted cut
generators and a valid implication chain to its conclusion. All these
certificates can be dovetailed. No negative reachability test, soundness
test, or computable rank is used.

For any fixed recorded formula \(\eta\), label \(\nu\), owner \(a\),
assignment \(s\), and variable \(x\), the numerical column \[
\{n:\Gamma_a^{<\nu}\vdash\eta^{s[x:=n]}\}
\] is therefore r.e. The finitely many values of \(s\) used by \(\eta\)
are hard-coded, and \(n\mapsto\eta^{s[x:=n]}\) is a computable syntax
operation.

\section{9. Erasure and the finite local prefix
lemma}\label{erasure-and-the-finite-local-prefix-lemma}

\subsection{Lemma 3: source sentence
necessity}\label{lemma-3-source-sentence-necessity}

For every source holder \(a\), \(S_a\vdash\chi\) implies
\(S_a\vdash K_a\chi\) for closed \(\chi\). This is the purely syntactic
necessity lemma proved in Section 3.1; it applies to the fixed theories
of Section 3 without a source-soundness assumption.

\subsection{Lemma 4: erasure}\label{lemma-4-erasure}

If \(\Gamma_a^{<\nu}\vdash\eta\) and \(\eta\) is closed, then
\(S_a\vdash\epsilon\eta\). The same assertion holds without a cut.

\textbf{Proof.} Each background row erases to a source theorem: modal PA
and structural rows erase to their source instances; assigned validity
erases to an assigned source validity by Lemma 1; reflection erases to
source truthfulness. A contextual-collapse row erases to
\(\operatorname{ucl}(\zeta\leftrightarrow\zeta)\), an actual validity.

A padded main row erases to an implication chain with the original
source code generator as conclusion. That generator is already a theorem
of \(S_a\), so propositional logic gives the padded implication,
irrespective of its marker antecedents. A generated own prefix erases to
an own prefix of a source theorem, and Lemma 3 applies.

Translate the finite valid implication of a cut proof by erasure using
Lemma 1. Its erased premises are all source theorems by the preceding
induction, so its erased conclusion follows. Numeric constants, source
formulas and assignments have not changed.

\subsection{Lemma 5: local prefixing of a derived cut
consequence}\label{lemma-5-local-prefixing-of-a-derived-cut-consequence}

Suppose \(a\in A\), \(\chi\) is closed and globally admissible,
\(\operatorname{supp}(\chi)<\nu<\delta\), and
\(\Gamma_a^{<\nu}\vdash\chi\). Then \[
\Gamma_a^{<\delta}\vdash K_a^\nu\chi.
\]

\textbf{Proof.} Choose a finite certificate with admitted cut premises
\(\eta_1,\ldots,\eta_m\). They are closed and globally admissible, and
their whole supports lie below \(\nu\), including any markers or earlier
prefix labels. The exact valid chain \[
V=\eta_1\to\cdots\to\eta_m\to\chi
\] is consequently closed, globally admissible, and supported below
\(\nu\). Its validity is the actual validity certificate, not a
conclusion drawn from erasure. Own validity admits \(K_a^\nu V\) below
\(\delta\). Own-prefix closure admits each \(K_a^\nu\eta_j\) below
\(\delta\). Same-label deduction distributes \(K_a^\nu\) through the
chain, yielding \(K_a^\nu\chi\). All new labels are \(\nu\), hence below
\(\delta\).

This proof uses arbitrary globally admissible background rows. The
premise list may contain different source-shared main-row tables, and
its combined implication need not possess one common source table.

\subsection{Lemma 6: markers can be discharged in their own
cut}\label{lemma-6-markers-can-be-discharged-in-their-own-cut}

Every marker \(K_a^h\tau\) with \(a\in A\) is derivable in every cut
strictly above \(h\). Therefore a padded main row whose whole support
lies below \(\delta\) yields its unpadded main sentence in
\(\Gamma_a^{<\delta}\).

\textbf{Proof.} The marker itself is an own-validity row, because
\(\tau\) is an actual closed arithmetic validity with empty support.
Repeated modus ponens discharges the marker antecedents of (M).

\section{10. Syntactic contextual equivalence and cut
collapse}\label{syntactic-contextual-equivalence-and-cut-collapse}

\subsection{Lemma 7: contextual substitution inside one
cut}\label{lemma-7-contextual-substitution-inside-one-cut}

If \(\mu\leq_1\nu\), the two fillings \[
\eta=C[K_b^\mu\phi],\qquad
\zeta=C[K_b^\nu\phi]
\] are globally admissible, and both supports lie below \(\delta\), then
for every assignment \(s\), \[
\Gamma_a^{<\delta}\vdash\eta^s
\iff
\Gamma_a^{<\delta}\vdash\zeta^s.
\tag{1}
\]

\textbf{Proof.} Numerically close the context, preserving the binders on
the path to the hole. More precisely, start with the partial
substitution assigning \(s(x)\) to every free variable. On entering a
quantifier binding \(x\), remove \(x\) from that substitution; transform
each fixed side formula using the current substitution. At the hole use
the resulting substitution on \(\phi\). This produces a context \(C_s\)
and a filler \(\phi_s\) satisfying the literal identities \[
\eta^s=C_s[K_b^\mu\phi_s],\qquad
\zeta^s=C_s[K_b^\nu\phi_s].
\] Both completed fillings are closed. They remain globally admissible
and have the same supports as before. Apply the contextual-collapse
admission rule to \(C_s,\phi_s,\mu,\nu\). Its universal closure is
vacuous, so the admitted formula is exactly
\(\eta^s\leftrightarrow\zeta^s\), and it lies below \(\delta\).
Propositional logic gives (1).

This proof deliberately does not infer numerical instances through modal
syntax from an arbitrary universal closure. The closed equivalence is a
new literal instance of the same contextual schema.

\subsection{Lemma 8: finite-pattern preservation of
generators}\label{lemma-8-finite-pattern-preservation-of-generators}

Let \(F\subseteq L\) be finite and \(f:L\to L\) preserve
\(<,\leq_1,\leq_2\) on \(F\). If \(\sigma\in\Gamma_a\) and
\(\operatorname{supp}(\sigma)\subseteq F\), then its label-renaming
\(f\sigma\) belongs to \(\Gamma_a\).

\textbf{Proof.} Inspect the admission certificate.

\begin{itemize}
\tightlist
\item
  Allowed owners do not change. Every nesting inequality is preserved
  because its two endpoints occur in the formula support.
\item
  Modal PA commutes with label renaming. Actual validity is preserved by
  Lemma 1, and assigned rows retain their assignment.
\item
  Structural and truthfulness rows retain their syntax. Their label
  inequalities remain true.
\item
  Both endpoints of a contextual-collapse row occur in its whole
  support, so its \(\leq_1\) premise remains true. Rename the arbitrary
  context and filler.
\item
  A main row retains its source code sentence, source numeral, owner,
  hereditary mask, tag formulas, and tag order. Replace its table by
  \(f\circ\ell\). Every displayed guard point occurs among its markers,
  so it belongs to \(F\). Renaming therefore preserves all guard-chain
  relations and all whole-argument \(\leq_1\) conditions. Freshness
  refers to unchanged source shapes and unchanged tags.
\item
  A prefix row remains a prefix row because its complete body and prefix
  label occur in the support, and every required strict bound is
  preserved.
\end{itemize}

Induction on the finite prefix certificate completes the proof. Forward
relation preservation suffices for this argument.

\subsection{Lemma 9: exact cut
collapse}\label{lemma-9-exact-cut-collapse}

If \(\alpha\leq_1\beta\) and \(\operatorname{supp}(\eta)<\alpha\), then
\[
\Gamma_a^{<\alpha}\vdash\eta
\iff
\Gamma_a^{<\beta}\vdash\eta.
\tag{2}
\]

\textbf{Proof.} The forward implication is monotonicity of cut
generators. For the reverse, take a finite certificate below \(\beta\),
and let \(F\) be the union of the supports of its conclusion and
displayed premises. Every point of \(F\) is below \(\beta\). In
particular, all guard markers and earlier prefix labels are included.

Finite reflection at \(\alpha\leq_1\beta\) gives a pattern map \(f\)
taking \(F\) below \(\alpha\) and fixing its points already below
\(\alpha\). Lemma 8 transports each displayed generator premise into the
cut below \(\alpha\). Lemma 1 transports the valid implication under
arbitrary label renaming. The conclusion is fixed because its support
was below \(\alpha\), so the reflected certificate proves \(\eta\).

The support \(F\) concerns the displayed premises and conclusion. The
logical validity derivation can be renamed globally regardless of the
labels in its intermediate lines; those intermediate lines need no
ordinal-pattern conditions. Thus one does not incorrectly demand that
arbitrary labels inside a validity derivation already lie below the
original cut.

\section{11. The filtered stage model and background truth before stage
soundness}\label{the-filtered-stage-model-and-background-truth-before-stage-soundness}

Let \(P=(P_a)_{a\in I}\) be alphabetic raw predicates on closed source
formulas. Its filtered source model \(F(P)\) uses \[
F(P)\models K_a\phi[s]
\iff P_a(\phi^s)\ \land\ F(P)\models\phi[s].
\] No soundness of \(P_a\) is assumed.

Define the stage model \(M=M(P)\) by standard arithmetic and the
following recursive modal clauses: \[
\begin{aligned}
M\models K_d\eta[s]
&\iff P_d((\epsilon\eta)^s)\land M\models\eta[s],\\
M\models K_a^\nu\eta[s]
&\iff \Gamma_a^{<\nu}\vdash\eta^s\land M\models\eta[s],\\
M\models Z_e\eta[s]&\iff\text{false}.
\end{aligned}
\tag{3}
\] These clauses apply to arbitrary recorded formulas, not merely
source-shared lifts.

\subsection{Lemma 10: the stage model is a legitimate base-logical
structure}\label{lemma-10-the-stage-model-is-a-legitimate-base-logical-structure}

The raw stage predicates are alphabetic. The filtered clauses satisfy
the modal structure conditions and the numerical-assignment law \[
M\models\eta[s]\iff M\models\eta^s[t]
\] for every \(s,t\).

\textbf{Proof.} Erasure takes alphabetic variants to alphabetic
variants, so live predicates inherit the property from \(P\). Closed
alphabetic variants are base-logically equivalent by the alphabetic
axioms; consequence closure therefore makes each cut predicate
alphabetic. Ghost predicates are empty.

Free-variable invariance, alphabetic invariance, and the legal
variable-substitution law now follow by formula induction. At a modal
head the raw test is applied to the numerically closed argument, and the
relevant closed arguments are identical or alphabetic variants by the
syntactic substitution laws; the truth conjunct follows from the
induction hypothesis. Numerical closure obeys the same recursive
identities and removes every free variable, giving the displayed
assignment law.

This constructs \(M\) before any stage-cut soundness theorem. Thus
actual validities are true in \(M\).

\subsection{Lemma 11: contextual collapse is true before stage
soundness}\label{lemma-11-contextual-collapse-is-true-before-stage-soundness}

For any admitted instance (CC), \(M\) satisfies its two fillings
equivalently at every assignment.

\textbf{Proof.} Induct on the one-hole context \(C\).

At the hole, global admissibility puts the support of \(\phi\) below
\(\mu\). Equation (2) equates the two raw proof tests
\(\Gamma_b^{<\mu}\vdash\phi^s\) and \(\Gamma_b^{<\nu}\vdash\phi^s\).
Their filtered truth conjunct is the same assertion \(M\models\phi[s]\).

Boolean contexts preserve equivalence, and quantified contexts use the
induction hypothesis at every updated assignment. At an unindexed outer
head the two argument erasures are literally equal, so the raw \(P_d\)
tests coincide; the inner truth conjuncts coincide by induction. At a
ghost outer head both values are false.

At an indexed outer head \(K_d^\xi D[-]\), both filled argument supports
lie below \(\xi\). Lemma 7, applied in holder \(d\)'s cut below \(\xi\),
equates the two raw cut tests. The induction hypothesis equates their
truth conjuncts. This use of Lemma 7 is syntactic membership of an
already admitted contextual row, not its assumed truth. Hence the
induction is noncircular.

Taking all assignments verifies the universal closure in (CC).

\subsection{\texorpdfstring{Lemma 12: every background row is true in
\(M\)}{Lemma 12: every background row is true in M}}\label{lemma-12-every-background-row-is-true-in-m}

\textbf{Proof.} Arithmetic equations hold in standard arithmetic. For
induction, the truth set \(\{n:M\models\eta[s[x:=n]]\}\) satisfies
ordinary induction on \(\mathbb N\). Its base and successor instances
are the literal schema instances because \(S(\bar n)=\overline{n+1}\)
and numerical closure commutes with the displayed \(0\)- and
successor-substitutions. This uses no unrestricted substitutivity of
extensionally equal arbitrary terms inside modalities.

An assigned validity is true by actual validity in \(M\) and Lemma 10's
assignment law. For own validity \(K_a^\nu\eta\), the argument is valid
and hence true. Its numerical closure \(\eta^s\) is an assigned-validity
row below \(\nu\), since its whole support lies below \(\nu\); this
supplies the raw proof test.

Same-label deduction follows from modus ponens in \(\Gamma_a^{<\nu}\),
together with the truth implications in the antecedent filters. For
introspection, an acceptance of \(K_a^\mu\eta\) gives both
\(\Gamma_a^{<\mu}\vdash\eta^s\) and truth of that modal atom. Lemma 5,
with \(\mu<\nu\), gives \(\Gamma_a^{<\nu}\vdash K_a^\mu\eta^s\), the
needed raw test for the outer head. Its truth conjunct is the accepted
inner modal atom.

Every truthfulness row follows directly from the filter in (3).
Contextual collapse follows from Lemma 11.

Each marker is also true before stage soundness: \(\tau\) is true, and
its empty-support validity is provable in \(\Gamma_a^{<h}\), so
\(K_a^h\tau\) is true. Therefore a padded main sentence is true
precisely when its unpadded code sentence is true.

\section{12. Coherent changes of labels and fixed-label
transport}\label{coherent-changes-of-labels-and-fixed-label-transport}

\subsection{Lemma 13: coherent lifts are interchangeable inside a
cut}\label{lemma-13-coherent-lifts-are-interchangeable-inside-a-cut}

Let \(t_1\theta,t_2\theta\) be two globally admissible lifts of the same
source formula, with the same hereditary live/ghost mask. Assume their
labels lie in a set \(C\) such that \[
u,v\in C,\ u\leq v\quad\Longrightarrow\quad u\leq_1v,
\] and both supports lie below \(\delta\). Then, for every active holder
\(a\) and assignment \(s\), \[
\Gamma_a^{<\delta}\vdash(t_1\theta)^s
\quad\Longleftrightarrow\quad
\Gamma_a^{<\delta}\vdash(t_2\theta)^s.
\]

\textbf{Proof.} Form the pointwise minimum table
\(\ell_0=\min(\ell_1,\ell_2)\) on each source owner/shape. It is
globally admissible: if both original child labels are below their
corresponding parent labels, choose a coordinate realizing the minimum
parent; the minimum child is at most the child in that coordinate and is
therefore strictly below the minimum parent.

Lower \(t_1\theta\) to the minimum lift occurrence by occurrence,
children before parents. At each step the changed label moves from \(v\)
to \(u\leq v\), with \(u\leq_1v\) by coherence. Its children have
already been lowered and lie below \(u\). Every unchanged indexed
ancestor is still above the old label \(v\), and hence above \(u\). All
other labels are unchanged. Thus both intermediate contexts are globally
admissible and stay below \(\delta\). Lemma 7 gives a cut equivalence
for this one occurrence. Compose finitely many such equivalences.

Do the same independently for \(t_2\theta\), then reverse its
equivalence to the minimum lift. Temporary intermediate formulas can
violate source sharing, but this does not violate any
contextual-collapse admission condition.

\subsection{Lemma 14: guarded fixed-label
transport}\label{lemma-14-guarded-fixed-label-transport}

Suppose \(a\in A\), \(G_{k+1}(\lambda)\), \(r(\beta)\leq k+1\), and a
prescribed globally admissible lift \(\beta'=t_{H,\ell}\beta\) keeps
every external query of \(\beta\) live. Assume \[
\xi\in\operatorname{supp}(\beta')
\quad\Longrightarrow\quad \xi<\lambda\ \text{and}\ \xi\leq_1\lambda.
\] Then for every assignment \(s\), \[
\Gamma_a^{<\lambda}\vdash(\beta')^s
\iff S_a\vdash\beta^s.
\tag{4}
\]

\textbf{Proof.} Erasure gives the forward implication. For reverse
transport one may replace \(H\) by \(Q(\beta)\), since the two masks
agree at every external query of \(\beta\), and thus give the same
target lift.

Choose a finite source certificate for \(\beta^s\), with premises
\(K_a^{n_j}\sigma_j\), where each \(\sigma_j\) is a source background or
code generator. Apply the mask \(Q(\beta)\) uniformly to the entire
certificate.

Define active depth \(D\) by counting only active modal heads, taking
maxima at Boolean branches, and passing through quantifiers and
nonactive heads. An active head over an argument of depth \(d\) will
receive a canonical label \(c_d\). Let \[
N=\sum_j D(K_a^{n_j}\sigma_j)+D(\beta)+1.
\] Every premise and the target have depth below \(N\).

Apply finite room at the guarded top \(\lambda\), above the finite
prescribed support. Obtain \(N\) separated copies
\(Y_0,\ldots,Y_{N-1}\), each containing a rank-\(k\) guard, with every
point below and \(\leq_1\)-below \(\lambda\). Let \(c_d=\min Y_d\).
Separation gives \(c_d<c_e\) for \(d<e\); connectedness gives
\(c_d\leq_1c_e\) for \(d\leq e\). Every original goal point and every
new copy point belong to the same \(\leq_1\)-coherent set below
\(\lambda\).

The canonical depth table is globally admissible on the whole finite
certificate. An active descendant query inside an argument of depth
\(d\) has strictly smaller depth and hence a smaller label. We now
derive each canonical recorded premise.

\textbf{Background generators.} Modal PA is carried to a literal PA
instance, because recording commutes with the substitutions in
induction. An assigned validity is carried to the same assignment of an
actually valid recorded formula by Lemma 1. Own validity is a listed
recorded own-validity row. Truthfulness becomes truthfulness of the
selected indexed, live, or ghost head. Introspection has two successive
active-depth labels and is its listed recorded instance.

Deduction requires an additional contextual-collapse step. Write \[
\gamma=c_{D(\phi\to\psi)},\quad
\mu=c_{D(\phi)},\quad \nu=c_{D(\psi)}.
\] The canonical image is \[
\operatorname{ucl}\bigl(
K_a^\gamma(t\phi\to t\psi)\to
(K_a^\mu t\phi\to K_a^\nu t\psi)\bigr).
\] Since \(D(\phi\to\psi)=\max(D\phi,D\psi)\), we have
\(\mu\leq_1\gamma\) and \(\nu\leq_1\gamma\). Begin with the admitted
same-label deduction row at \(\gamma\). Change the displayed occurrence
over \(t\phi\) from \(\gamma\) to \(\mu\), and the displayed occurrence
over \(t\psi\) from \(\gamma\) to \(\nu\), using two contextual-collapse
equivalences. Both final and intermediate formulas are admissible: the
relevant arguments were already below their own smaller depth labels.
All introduced labels were already in the canonical mixed row, so no
cutoff increases. The universal-closure binders can be included in the
contexts; no numerical specialization is needed.

\textbf{Code generators.} Let \(\sigma_j=C_{w,\phi}\). The target \(w\)
is passive and nonexternal by the side conditions, so the canonical
image has the correct literal main-code shape and unchanged numeral. A
live external argument \(\chi\) in this image survives only because its
owner/shape occurs in \(Q(\beta)\). Thus \[
r(\chi)<r(\beta)\leq k+1,\qquad r(\chi)\leq k.
\] An active occurrence at depth \(d\) inside \(\chi\) receives the
first \(r(\chi)+1\) points of \(Y_d\) as its guard. Every active label
inside its own argument has smaller depth; connectedness supplies the
required \(\leq_1\) relation to \(c_d\).

Let \(D_j=D(\sigma_j)\). All guards needed by this main row can be taken
inside \(\bigcup_{d<D_j}Y_d\). Choose fresh marker tags for these finite
points. If the source premise has no own prefix, every marker and every
main label lies below \(\lambda\), so Lemma 6 derives the canonical
unpadded code below \(\lambda\). If it has at least one own prefix,
separation puts every point of \(\bigcup_{d<D_j}Y_d\) below \(c_{D_j}\),
the first own-prefix label; derive the base code already in that smaller
cut.

\textbf{Own prefixes.} The equality \[
D(K_a^n\sigma)=D(\sigma)+n
\] holds because \(a\in A\). Starting from the base proof below
\(c_{D(\sigma)}\), repeatedly use Lemma 5 at the successive labels \[
c_{D(\sigma)},c_{D(\sigma)+1},\ldots,
c_{D(\sigma)+n-1}.
\] Every body, including the guards and marker proofs used to derive it,
is in the preceding cut. The final prefix label is below \(\lambda\).
This derives the full recorded source premise in
\(\Gamma_a^{<\lambda}\).

The uniformly recorded finite source implication is actually valid by
Lemma 1. Combining the recorded premise proofs derives the canonical
recorded target below \(\lambda\). Lemma 13 then identifies it with the
prescribed target because all their labels lie in the common coherent
set and both lifts are admissible.

In fact, room placed every canonical label above every prescribed goal
label, so the final step can use direct child-before-parent lowering
rather than pointwise minima. This is a simplification, not an extra
assumption.

The number \(N\), the source certificate, and the reflected or auxiliary
proof labels may depend on \(s\). The guarded top, target table and
target shape do not. Formula variants of the same source shape have the
same rank, support and admissibility conditions, so the argument applies
separately to every such variant. No uniform numerical proof-depth bound
has been used.

\section{13. Exact recovery from one finite
room}\label{exact-recovery-from-one-finite-room}

\subsection{Lemma 15: recovery at room-copy
heads}\label{lemma-15-recovery-at-room-copy-heads}

Let \(Y_0,\ldots,Y_{N-1}\) be room copies below a top \(\lambda\) with
guard \(G_{k+1}(\lambda)\), and let \(h_i=\min Y_i\). Suppose a
ghost-free prescribed lift \(t\beta\) is admissible,
\(r(\beta)\leq k+1\), and every target label is some \(h_d\) with
\(d<i<N\). Then \[
\Gamma_a^{<h_i}\vdash(t\beta)^s
\iff S_a\vdash\beta^s
\] for every active holder \(a\) and assignment \(s\).

\textbf{Proof.} Erasure is forward. Lemma 14 proves the target below
\(\lambda\). Room gives \(h_i\leq_1\lambda\), while separation places
the target support below \(h_i\). Lemma 9 collapses the finite proof to
\(h_i\), fixing every target label.

For a given source formula \(\theta\), choose \(k=r(\theta)\),
\(N=D(\theta)+1\), and one such room using finite tallness and room.
Record every active head over an argument of depth \(d\) at \(h_d\), and
keep every external query of \(\theta\) live. Every active query
\((a,\operatorname{sh}\psi)\in Q(\theta)\) has \(D(\psi)<D(\theta)\),
and all labels inside its recorded argument have smaller indices. Lemma
15 gives \[
\Gamma_a^{<h_{D(\psi)}}\vdash t\psi
\iff S_a\vdash\psi
\tag{5}
\] for every closed formula \(\psi\) of that query shape. On nonclosed
formulas both predicates are false. Thus (5) covers every term variant
and every numerical instance in the class, not only the literal
occurrence in \(\theta\).

If the original active predicates are full source theoremhood, (5),
literal live erasure, and formula induction give \[
M(P)\models t\theta[s]\iff F(P)\models\theta[s].
\tag{6}
\] Only the finitely many query shapes of this \(\theta\) are required.

\section{14. Reconstruction with full retained
predicates}\label{reconstruction-with-full-retained-predicates}

The stage language is a temporary instrument. To apply a lower-ideal
genericity hypothesis, we need a source-language model in which the full
lower theories are still present. The following reconstruction retains
those predicates in full, uses stage cuts only for the active query
classes of the target, and keeps the required external queries. Its
soundness conclusion covers every accepted sentence of these predicates.

The following statement is the precise form needed in the lower-rank
argument.

Fix a source formula \(\theta\), a hereditary mask \(H\), and an
admissible recorded lift \(t\theta\). Let
\(D_0\subseteq I\setminus(A\cup E)\) be the retained set. It can be
infinite. Assume every nonactive, nonexternal owner appearing in
\(Q(\theta)\) either belongs to \(D_0\) or has empty old predicate
\(P_d\).

Define a new original-language raw family \(Q'\) as follows, on closed
formulas \(\psi\): \[
Q'_d(\psi)=
\begin{cases}
P_d(\psi),&d\in D_0,\\
\bigl((d,\operatorname{sh}\psi)\in Q(\theta)\bigr)
\land\bigl(\Gamma_d^{<\ell(d,\operatorname{sh}\psi)}\vdash t\psi\bigr),
&d\in A,\\
\bigl((d,\operatorname{sh}\psi)\in Q(\theta)\cap H\bigr)\land P_d(\psi),
&d\in E,\\
\text{false},&\text{otherwise}.
\end{cases}
\tag{7}
\] All nonclosed formulas are rejected. Full retained predicates in the
first branch are not restricted by shape and are not recovered from
indexed cuts.

\subsection{Lemma 16: reconstruction truth and column
admissibility}\label{lemma-16-reconstruction-truth-and-column-admissibility}

For every formula \(\psi\) with \(Q(\psi)\subseteq Q(\theta)\) and every
assignment \(s\), \[
F(Q')\models\psi[s]\iff M(P)\models t\psi[s].
\tag{8}
\] If \(P\) is column admissible, so is \(Q'\).

\textbf{Proof of truth.} At an active modal occurrence the two raw tests
agree by the second branch of (7) and commutation of recording with
numerical closure. At a retained passive owner, erasure
\(t\chi\mapsto\chi\) makes the live raw test exactly \(P_d(\chi^s)\). At
a live external owner the third branch retains precisely this raw test;
at a ghost it is false on both sides. At any remaining queried passive
owner the hypothesis makes its old predicate empty, matching the final
branch. The truth filters then agree by induction on the argument.
Boolean connectives preserve equivalence, and quantifiers use the same
induction hypothesis for every updated assignment.

\textbf{Proof of admissibility.} Alphabetic variation preserves shapes
and commutes with recording; old predicates are alphabetic, and cut
predicates are alphabetic by Lemma 10. Every branch in (7) is therefore
alphabetic.

Fix one owner \(d\), template \(\chi\), variable \(x\), and parameter
assignment \(s\). All \(\chi^{s[x:=n]}\) have the same shape, so the
choice among branches of (7) is independent of \(n\). The column is an
old r.e. column, an empty column, or a column of one fixed recorded cut
and one fixed finite recorded template \(t\chi\). Lemma 2 supplies the
last column. Hence every numerical parameter column is r.e. No jointly
computable choice of \(\ell\), \(D_0\), or program indices has been
asserted or needed.

\subsection{Lemma 17: full soundness of the new outside
predicates}\label{lemma-17-full-soundness-of-the-new-outside-predicates}

In addition to the preceding hypotheses, assume:

\begin{itemize}
\tightlist
\item
  \textbf{S:} For every \((a,q)\in Q(\theta)\) with \(a\in A\), the raw
  predicate \(\Gamma_a^{<\ell(a,q)}\vdash-\) is sound in \(M(P)\) on
  every closed recorded formula, including nonadmissible conclusions.
\item
  \textbf{E:} Every old \(P_e\), \(e\in E\), is fully sound in \(F(P)\).
\item
  \textbf{X:} For every \((e,\operatorname{sh}\chi)\in Q(\theta)\cap H\)
  with \(e\in E\), every \((a,\operatorname{sh}\beta)\in Q(\chi)\) with
  \(a\in A\), and every closed source formula \(\psi\) with
  \(\operatorname{sh}\psi=\operatorname{sh}\beta\), \[
  \Gamma_a^{<\ell(a,\operatorname{sh}\beta)}\vdash t\psi
  \quad\Longleftrightarrow\quad P_a(\psi).
  \] The query set \(Q(\chi)\) depends only on its shape, so any
  representative \(\chi\) of that shape gives the same condition.
\end{itemize}

Then every \(Q'_d\) with \(d\notin D_0\) is fully sound in \(F(Q')\).

\textbf{Proof.} A retained active acceptance \(Q'_a(\psi^s)\) gives a
cut proof of \((t\psi)^s\). The corresponding label occurs in
\(t\theta\) because the owner/shape class was retained. S gives
\(M(P)\models t\psi[s]\); (8) gives \(F(Q')\models\psi[s]\). This
argument applies to every accepted formula of that shape.

For an external acceptance, \((e,\operatorname{sh}\psi)\) lies in both
\(Q(\theta)\) and \(H\), and \(P_e(\psi^s)\) holds. E gives
\(F(P)\models\psi[s]\). Heredity makes every external head inside
\(t\psi\) live. At every nested active query, X identifies the new raw
test with \(P_a\); at retained passive heads the old predicates are
literal; at nested external heads the third branch retains the old test;
every other queried passive predicate is empty on both sides. Thus
\(Q'\) and \(P\) agree on every query shape of \(\psi\), including all
variants and numerical instances. Formula induction transfers its old
truth to \(F(Q')\).

All other outside predicates are empty. This proves full outside
soundness. No predicate in \(D_0\) was assumed sound.

\subsection{How literal main guards supply
X}\label{how-literal-main-guards-supply-x}

At a live external occurrence \(K_e\chi\), the main-row guard rule gives
\(G_{r(\epsilon\chi)}(\nu)\) at every indexed descendant
\(K_b^\nu\beta\), and puts the entire support of \(\beta\) strictly and
\(\leq_1\)-below \(\nu\). Its source argument shape has rank strictly
below \(r(\epsilon\chi)\). The retained external subtree is ghost-free
by heredity. Apply Lemma 14 at \(\nu\), with \(k+1=r(\epsilon\chi)\), to
each formula of the descendant's source shape. Erasure supplies the
reverse direction. If \(P_b=\operatorname{Th}(S_b)\) on active owners,
this is exactly X.

The hypotheses are shape invariants, so the lemma holds for the full
infinite class of term variants rather than only the inspected literal
occurrence.

\section{15. Local main truth and the two independent
inductions}\label{local-main-truth-and-the-two-independent-inductions}

The induction on the given hierarchy supplies soundness for strictly
lower source agents. Inside one step of that induction, a second
induction on ordinal support proves truth of the stage generators.
Main-row truth uses the already sound smaller cuts together with a
strictly lower hierarchy frontier. This section establishes that inner
step under the stated outer hypotheses; Section 16 then discharges the
outer hypotheses for every frontier.

The following argument supplies the inner induction step used in the
proof of \(G(J)\) in Section 16. Its only outer induction hypotheses are
\(G(J')\) for strictly lower frontiers, and the direct core lemma of
Section 5.

Fix a finite frontier \(J=A\), let \(P\) be column-admissible and retain
the full core and all source theories in the downward ideal of \(J\),
and suppose outside predicates are sound, with only the finite \(E\)
possibly nonempty. Suppose genericity has already been proved for every
strictly lower nonempty finite frontier, and the direct auxiliary-core
genericity lemma is available.

Consider a main row for active holder \(a\) with target \(w=*\) or
\(wRa\), source argument \(\phi\), and table \(t\). Collect \(w\) when
original, together with all original passive owners occurring in
\(\phi\) that lie strictly below \(J\). Take the maximal finite frontier
\(J'\) of this finite set, and let \[
D_0=\{*\}\cup\downarrow J'.
\] Here \(\downarrow J'=D_R(J')\) denotes the entire lower ideal,
including the frontier itself. Every member of \(J'\) lies strictly
below some member of \(J\). If \(J'\ne\varnothing\), its maximum
well-founded \(R\)-rank is strictly below that of \(J\). If the finite
set is empty, \(w=*\) and core genericity is the appropriate base case.

Every original passive full owner named in \(\phi\), and \(w\) itself,
lies in \(D_0\). No member of \(D_0\) is active or external, by
antichainness and the choice of \(E\). An old outside owner outside
\(E\) is empty. These facts verify exactly the passive-owner condition
for (7).

Assume all cuts at labels in the main row's support are already sound.
These are S. Old external soundness is E. Literal guards supply X by the
preceding paragraph. Lemmas 16--17 construct a column-admissible \(Q'\),
retain the entire \(D_0\) predicates, and make every outside predicate
fully sound. Lower-rank genericity, or core genericity when \(J'\) is
empty, therefore makes \(S_w\) sound in \(F(Q')\).

For each \(n\) with \(S_w\vdash\phi(\bar n)\), equation (8) now gives
\(M(P)\models t\phi[x:=n]\). The outer live head \(K_w\) has raw
predicate exactly \(S_w\)'s theoremhood, and the arithmetic side still
uses the actual source numeral \(p(w,\phi)\). Consequently both sides of
\[
\forall x\bigl(K_wt\phi\leftrightarrow
W(\overline{p(w,\phi)},x)\bigr)
\] coincide for every \(n\): forward uses its raw proof test and the
source-column equation; reverse uses that equation and the just-proved
recorded-body truth. The markers are independently true by Lemma 12.
This proves the padded main row.

\subsection{Lemma 18: support induction yields soundness of every stage
cut}\label{lemma-18-support-induction-yields-soundness-of-every-stage-cut}

Every \(\Gamma_a\)-generator is true in \(M(P)\), and every cut
predicate is fully sound there.

\textbf{Proof.} Prove generator truth simultaneously for all active
holders by well-founded induction on the maximum of nonempty finite
supports, with an initial empty-support case.

At empty support, there are no indexed queries or prefix labels.
Background rows are true by Lemma 12. A main row is covered by the local
argument above, with S and X vacuous where no indexed labels occur. No
prefix row has empty support.

For a label \(\nu\), assume every generator whose whole support is
strictly below \(\nu\) is true. Every cut at a label \(\mu\leq\nu\) is
then sound: every generator below \(\mu\) has support strictly below
\(\nu\), and base-logical soundness transfers their truth to all finite
cut consequences.

Take a generator with maximum support \(\nu\). A background row is
already true. A main row is true by the local reconstruction argument,
since cuts at every label in its support are now sound. A prefix
\(K_a^\mu\sigma\) has \(\mu=\nu\) because its body support lies strictly
below \(\mu\). The induction hypothesis makes \(\sigma\) true; its
literal generator membership below \(\mu\) supplies the raw cut test.
The prefix is therefore true.

This proves every finite-support generator. Base-logical soundness now
makes every cut fully sound on all its closed consequences, including
consequences that are not globally admissible. No stage proof was used
to assume the original family's soundness. The outer original-rank
induction was used only for the strictly lower \(J'\) selected in the
main-row argument.

\section{16. Whole-ideal genericity and the actual
family}\label{whole-ideal-genericity-and-the-actual-family}

We now discharge the outer induction hypotheses used in Section 15.

\subsection{16.1. The outer induction}\label{the-outer-induction}

Since \(R\) is well founded, let \[
\rho(i)=\sup_{jRi}(\rho(j)+1).
\] Then \(jR^+i\) implies \(\rho(j)<\rho(i)\). For a nonempty finite
antichain \(J\), put \[
r(J)=\max_{j\in J}\rho(j).
\] This is a metatheoretic ordinal rank. It is never queried by an
enumerator.

Prove \(G(J)\), as defined in Section 4.3, by induction on \(r(J)\) for
nonempty \(J\). The empty frontier is the separately proved core lemma
\(G(\varnothing)\). In particular, no assertion about a maximum of the
empty set is used.

Fix a nonempty \(J\) and assume \(G(J')\) for every nonempty finite
frontier of smaller rank.

\subsection{16.2. Reduction to finitely many nonempty outside
owners}\label{reduction-to-finitely-many-nonempty-outside-owners}

It suffices to prove \(G(J)\) for families whose nonempty original
outside predicates have their owners in a finite set \(E\).

Indeed, let \(P\) be an arbitrary family satisfying the hypotheses of
\(G(J)\), and let \(\theta\) be a theorem of a retained source owner.
Define \(\widehat P\) by retaining the entire core and ideal unchanged
and setting, for outside owners, \[
\widehat P_b(\psi)=
P_b(\psi)\ \land\
\bigl((b,\operatorname{sh}\psi)\in Q(\theta)\bigr)
\] on closed \(\psi\). This family is column-admissible by Section 4.2.
Its nonempty outside owners belong to the finite set of outside names
occurring in \(Q(\theta)\).

Every retained outside acceptance has an argument shape from
\(Q(\theta)\). All queries inside an argument of that shape also lie in
\(Q(\theta)\), by heredity. Thus \(F(P)\) and \(F(\widehat P)\) agree on
that argument. Old outside raw soundness transfers to \(\widehat P\),
making every outside predicate raw-sound in \(F(\widehat P)\).

The finite-outside version of \(G(J)\) makes \(\theta\) true in
\(F(\widehat P)\). Agreement on \(Q(\theta)\) makes it true in \(F(P)\).
Applying this argument separately to every retained theorem, including
core theorems, proves the full \(G(J)\). This semantic reduction does
not change the source family, its theories, or any program numeral.

We may therefore fix a finite \(E\subseteq\mathbb N\setminus D_R(J)\)
containing all nonempty original outside owners.

\subsection{16.3. Stage soundness at the current
frontier}\label{stage-soundness-at-the-current-frontier}

Set \(A=J\), and build the extension stages of Sections 7--15. Their
side conditions hold: the core is passive, \(E\) is external to the
ideal, and every direct subordinate of an active owner is passive and
nonexternal.

The local argument of Section 15 selects, from a main source-code
argument, a finite set of strictly lower owners, including its code
owner when that owner is original. Every member of its maximal frontier
\(J'\) lies strictly below some member of \(J\). Therefore, whenever
\(J'\ne\varnothing\), \[
r(J')<r(J).
\] The available outer induction hypothesis is exactly the lower
genericity required there. When \(J'=\varnothing\), Section 5 is
available instead. Lemma 18 now proves that every stage cut is raw-sound
in the stage model \(M(P)\).

The induction order is thus fixed. Syntactic transport, cut collapse,
and background truth were proved first. Main-row truth uses sound
smaller stage cuts and genericity of a strictly lower original frontier.
Neither operation assumes soundness of an active original source theory.

\subsection{16.4. Recovery of the entire original
ideal}\label{recovery-of-the-entire-original-ideal}

First take an active owner \(a\in J\) and an arbitrary closed theorem
\(\theta\) of \(T_a\). Choose the finite room for \(\theta\) from Lemma
15, and let \(t\theta\) be its fully live room lift.

Apply Lemma 14 at the room top \(\lambda\). It gives \[
\Gamma_a^{<\lambda}\vdash t\theta.
\] Stage-cut soundness makes \(t\theta\) true in \(M(P)\). Room recovery
at every active query shape in \(\theta\), together with literal erasure
at passive heads, identifies that truth with \(F(P)\models\theta\). Thus
every active source predicate is raw-sound in \(F(P)\).

Now let \(i\in D_R(J)\setminus J\) and take an arbitrary closed
\(T_i\)-theorem \(\theta\). Choose a fully live room lift \(t\theta\) as
before. Form the finite set consisting of \(i\) and every original owner
named in \(\theta\) that lies strictly below \(J\). Let \(J'\) be its
\(R^+\)-maximal frontier, and retain \[
D_0=\{*\}\cup D_R(J').
\] This frontier is nonempty, contains \(i\) in its downward ideal, and
has smaller rank than \(J\).

Apply the reconstruction of Section 14. Every full passive original
owner named in \(\theta\) belongs to \(D_0\); any other nonactive,
nonexternal queried owner has empty old predicate. All used stage cuts
are now sound, giving S. Old outside soundness gives E. Room recovery
gives X at all active query shapes of \(\theta\), and hence in
particular below every retained external argument. Lemmas 16--17 give a
column-admissible family \(Q'\) that keeps the full source predicates on
\(D_0\) and has fully sound predicates outside \(D_0\).

By \(G(J')\), the unchanged \(T_i\)-predicate is sound in \(F(Q')\), so
\(F(Q')\models\theta\). Reconstruction transfers this to
\(M(P)\models t\theta\), and room agreement transfers it to
\(F(P)\models\theta\). This proves raw soundness for every passive owner
in the retained ideal.

These arguments are performed separately for each theorem. The ideal may
be infinite, and unrelated theorems need not use related rooms or label
tables.

Every original raw predicate is now sound in \(F(P)\): those inside the
ideal by the preceding arguments, and those outside by hypothesis. Apply
the direct core lemma at this point to obtain raw soundness of \(U\).
This proves the finite-outside version of \(G(J)\), hence its full
version by Section 16.2. The outer induction is complete.

\subsection{16.5. Soundness of the actual source
family}\label{soundness-of-the-actual-source-family}

Let \(P^{\mathrm{all}}_a=\operatorname{Th}(S_a)\) for all \(a\in I\). It
is column-admissible. To prove original raw soundness, fix a closed
\(T_i\)-theorem \(\theta\).

Take the finite set consisting of \(i\) and every original modal owner
occurring in \(\theta\), and let \(J\) be its \(R^+\)-maximal frontier.
Retain the full source core and \(D_R(J)\), and empty every other
original raw predicate. Call the resulting family \(P^J\).

Empty outside predicates are raw-sound without any assumption about
\(P^{\mathrm{all}}\). Thus \(G(J)\) gives \(F(P^J)\models\theta\). Every
query in \(\theta\) has retained owner, so query agreement gives \[
F(P^{\mathrm{all}})\models\theta.
\] This proves raw soundness of every original source theory in the
actual filtered model. Apply core genericity once more to get raw
soundness of \(U\).

All truth filters can now be removed. We obtain a genuine
standard-arithmetic structure \(\mathcal M_*\) with \[
\mathcal M_*\models K_a\varphi[s]
\quad\Longleftrightarrow\quad
S_a\vdash\varphi^s,
\qquad a\in I,
\tag{A1}
\] in which every source theorem is true.

\section{17. The required schemas and the effective
reduct}\label{the-required-schemas-and-the-effective-reduct}

\subsection{17.1. Epistemic Arithmetic and known
truthfulness}\label{epistemic-arithmetic-and-known-truthfulness}

Every original theory contains its displayed background and all finite
own prefixes of that background. The source-family soundness theorem
consequently verifies modal PA, knowledge of its axioms, own validity,
deduction, introspection, their known versions, and reflection. These
include all axioms of self-doubting Epistemic Arithmetic.

For every \(i\in\mathbb N\), \(b\in I\), and unrestricted
enlarged-language formula \(\varphi\), the closed reflection sentence \[
\sigma=\operatorname{ucl}(K_b\varphi\to\varphi)
\] belongs to \(T_i\). By (A1), \(\mathcal M_*\models K_i\sigma\). This
proves full mutual known truthfulness in the enlarged language.

\subsection{17.2. Known particular subordinate
codes}\label{known-particular-subordinate-codes}

For \(jRi\) and unary \(\varphi\), let \(e=c(q,j,\varphi)\). The exact
closed sentence \[
C_{j,\varphi}(e)=
\forall x\bigl(K_j\varphi\leftrightarrow W(\bar e,x)\bigr)
\] is a \(T_i\)-generator. Thus \[
\mathcal M_*\models K_iC_{j,\varphi}(e).
\] Assigning the fresh variable \(e\) this numerical value witnesses
(Code). Under (A1), the raw modal test closes that variable to exactly
the source numeral appearing in the generator. This is an application of
the numerical-assignment semantics, with no modal term-generalization
principle.

The same argument gives each original agent a known particular code for
every unary core column.

\subsection{17.3. Entire-core
inheritance}\label{entire-core-inheritance}

\textbf{Inheritance lemma.} Every closed \(U\)-theorem is a theorem of
every \(T_i\).

\textbf{Proof.} Suppose \(U\vdash\sigma\), with \(\sigma\) closed, and
choose a variable \(x\). Put \[
\psi(x)=\sigma\land x=x.
\] This has exactly the one free variable \(x\), and
\(U\vdash\psi(x\mid\bar0)\). Equation (S4) gives \(0\in W_{d(\psi)}\). A
finite positive computation trace supplies \[
\mathrm{PA}\vdash W(\overline{d(\psi)},\bar0),
\] hence the same theorem in \(T_i\).

The core-code generator in \(T_i\) is \[
\forall x\bigl(K_*\psi\leftrightarrow
                   W(\overline{d(\psi)},x)\bigr).
\] Use the Assigned Validity numerical-instantiation argument of Section
2.2 to obtain its \(0\)-instance. The positive arithmetic fact gives
\(T_i\vdash K_*\psi(x\mid\bar0)\). The closed core-reflection instance
then gives \(T_i\vdash\psi(x\mid\bar0)\), hence \(T_i\vdash\sigma\).
\(\square\)

This is a proof-theoretic inheritance argument, using the literal
source-code numeral. It does not replace a code by one for a recorded or
masked argument.

\subsection{17.4. Known mechanicalness}\label{known-mechanicalness}

For any original \(j\), the sentence \(\operatorname{MC}_j(\varphi;x)\)
belongs to \(U\). Inheritance puts it in every \(T_i\), so (A1) gives \[
\mathcal M_*\models K_i\operatorname{MC}_j(\varphi;x).
\tag{A2}
\] When \(\operatorname{FV}(\varphi)=\{x\}\), this is precisely
(Mechanical). When there are other free parameters, (A2) gives the
stronger parameter schema recorded in Section 18.1.

\subsection{17.5. Removing the auxiliary
operator}\label{removing-the-auxiliary-operator}

For each original \(i\), define \[
T_i^0=
\{\sigma:\sigma\text{ is a closed }\mathcal L\text{-formula and }
                         T_i\vdash_{\mathcal L_*}\sigma\}.
\tag{A3}
\] This is uniformly r.e.: enumerate original source theorems and filter
by the decidable condition of containing no \(K_*\).

It is consequence-closed in the original base logic. Indeed, a finite
original-language proof from members of \(T_i^0\) has a validity
certificate that remains valid in the enlarged language, since every
enlarged-language structure has an original-language reduct. Its
premises are \(T_i\)-theorems, so its conclusion is a \(T_i\)-theorem
and therefore belongs to \(T_i^0\).

For every closed original-language sentence \(\sigma\), \[
T_i^0\vdash_{\mathcal L}\sigma
\quad\Longleftrightarrow\quad
T_i\vdash_{\mathcal L_*}\sigma.
\tag{A4}
\] The reverse direction holds by the definition of \(T_i^0\); the
forward direction was just proved.

Consequently, formula induction identifies the original-language reduct
of \(\mathcal M_*\) with the actual consequence interpretation of
\((T_i^0)\). All the required original-language formulas, their modal
induction instances, and their code sentences survive the reduct. In
particular, the original source indices still enumerate the exact
original unary columns by (A4).

Renaming \(T_i^0\) as \(T_i\) proves the theorem of Section 1.
\(\square\)

\section{18. Further consequences and
scope}\label{further-consequences-and-scope}

The same source family yields parameter mechanicalness and entire-theory
inheritance. We then examine a self-mechanical extension of the
auxiliary theory, construct a distinguished common agent, and specify
the quantifier order and the self-code boundary of the result.

\subsection{18.1. Known mechanicalness with
parameters}\label{known-mechanicalness-with-parameters}

For every original \(i,j\) and every original-language formula
\(\varphi(x,\vec z)\), where \(\vec z\) lists its other free variables,
the construction gives \[
\mathcal M\models
K_i\left[
\forall\vec z\ \exists e\ \forall x
 \bigl(K_j\varphi(x,\vec z)\leftrightarrow W(e,x)\bigr)
\right].
\tag{P1}
\] The universal parameter quantifiers are inside \(K_i\) and outside
\(\exists e\). The witness index may depend on the values of \(\vec z\).
This is a known universally quantified sentence, not merely a separate
known sentence for each externally fixed numeral tuple.

The proof is direct: (S3) includes that closed core generator, Section
17.3 transfers the entire sentence, and the reduct preserves it.

This is an additional explicit schema beyond the unary formulation of
the target. No claim of its logical independence from the exact unary
target plus the surrounding weak-modal arithmetic is needed or made.
Moving parameter quantifiers through knowledge would require its own
argument.

\subsection{18.2. Entire-theory inheritance and transitive subordinate
codes}\label{entire-theory-inheritance-and-transitive-subordinate-codes}

If \(jRi\), the inheritance proof of Section 17.3 works with \(T_j\) in
place of \(U\), using \(i\)'s particular \(j\)-codes. Hence, first in
the enlarged language and then in the reduct, \[
\operatorname{Th}(T_j)\subseteq\operatorname{Th}(T_i).
\tag{P2}
\] Iterating along a finite path gives the same inclusion whenever
\(jR^+i\).

A particular code sentence for a lower agent, proved at one step of a
path, is therefore inherited by every agent above that step. Thus the
family has known particular codes for all \(R^+\)-subordinates. The
theorem does not assert that \(R^+\) is exactly the relation of all
pairs for which such codes happen to be known; that would require
additional negative assertions about nonedges.

This inheritance is a consequence of the 3-1-3 requirements with the
standard \(W\)-representation and PA's positive computation
certificates. In a general solution, own validity and the actual
consequence semantics give each theory every Assigned Validity instance:
validity of \(\eta\) makes \(K_i\eta\) true at every assignment, hence
\(T_i\vdash\eta^s\). Known truthfulness makes the selected known code
sentence true; a lower theorem then gives a true concrete positive
\(W\)-fact, and the same padding and reflection argument transfers the
theorem. Alexander explicitly discusses the propagation to transitive
closure in {[}A, §2.6, p.~75{]}.

Own prefix closure gives all finite own iterations of known assertions.
Inheritance also gives appropriate nesting along subordinate chains: if
\(kR^+jR^+i\), \(\sigma\) is closed, and \(T_k\vdash\sigma\), then \[
\mathcal M\models K_iK_jK_k\sigma.
\] These observations do not establish arbitrary mixed-agent common
knowledge across unrelated agents.

\subsection{18.3. A strict self-mechanical extension of the auxiliary
theory}\label{a-strict-self-mechanical-extension-of-the-auxiliary-theory}

The auxiliary theory's generator set omits its own mechanicalness. To
measure the effect of adding those generators, write \(U^-=U\), and
define in the same enlarged language \[
U^+=\operatorname{cl}_*\left(
 B_*+\{\operatorname{MC}_b(\varphi;x):b\in I,\ \varphi,\ x\}
\right).
\tag{P3}
\] This theory adds the core's own mechanicalness instances to the
generator table of Section 3.2.

The inclusion \(U^-\subseteq U^+\) is immediate, but strictness requires
proof.

\textbf{Strictness proposition.} In the enlarged base logic, \[
\operatorname{Th}(U^-)\subsetneq\operatorname{Th}(U^+).
\]

\textbf{Proof.} Define an arithmetic translation \(E\) by \[
E(K_*\psi)=E(\psi),\qquad
E(K_j\psi)=\bot\quad(j\in\mathbb N),
\] and let it commute with arithmetic symbols, Boolean operations, and
quantifiers. Interpret formulas over standard \(\mathbb N\) by \[
M_E\models\varphi[s]\quad\Longleftrightarrow\quad
\mathbb N\models E(\varphi)[s].
\] The translation respects alphabetic changes and capture-free
substitution, and its truth depends only on free variables. Thus this is
an admissible base-logical structure with the numerical-assignment
property.

All \(B_*\)-generators are true. Modal induction is induction on a
subset of standard \(\mathbb N\). Core validity, deduction,
introspection, and core reflection become their transparent classical
counterparts. Original-owner reflection has a false antecedent. Actual
validity and Assigned Validity hold as in Section 2.

Every original knowledge column is empty, so a fixed empty-program index
witnesses every original-owner mechanicalness generator, including its
parameter version. A core prefix preserves truth because \(K_*\) is
transparent in this model. Hence \(M_E\models U^-\).

Let \[
\chi(x)=\neg W(x,x).
\] Its arithmetic truth set is the diagonal nonhalting set \[
D=\{n:n\notin W_n\}.
\] This set is not r.e.: if \(D=W_d\), then \(d\in W_d\) would be
equivalent to \(d\notin W_d\). In \(M_E\), the column for \(K_*\chi(x)\)
is exactly \(D\). Therefore \[
\operatorname{MC}_*(\chi;x)=
\exists e\,\forall x
 \bigl(K_*\chi(x)\leftrightarrow W(e,x)\bigr)
\tag{P4}
\] is false in \(M_E\).

By base-logical soundness, \(U^-\) does not prove (P4). The theory
\(U^+\) contains it as a generator. This proves strict inclusion.
\(\square\)

The separating structure is a base-logical countermodel, not a claim
that its transparent core is itself an r.e. knowing machine. Arbitrary
admissible structures are the correct models for this nonprovability
argument.

In the source model constructed in Section 16.5, since \(U^-\) is r.e.,
its own actual \(\chi\)-knowledge column is r.e., so (P4) is externally
true there. But \(U^-\) does not prove (P4), so the actual \(U^-\) core
does not know this instance of its own mechanicalness. The sentence is a
generator of \(U^+\).

If the source construction is carried out with \(U^-\) or \(U^+\) as its
auxiliary theory, denote the resulting original-agent theories by
\(T_i^-\) and \(T_i^+\), respectively. The strictness result concerns
the auxiliary theories; it does not prove
\(\operatorname{Th}(T_i^-)\subseteq\operatorname{Th}(T_i^+)\). Changing
the core changes the literal code numerals and the recursion fixed
point, so the original generator sets are not simply nested.

\subsection{18.4. A distinguished self-mechanical common
agent}\label{a-distinguished-self-mechanical-common-agent}

Given \(R\), add a new least agent \(*\), with \[
*\,R^*\,i\quad\text{for every original }i,
\] retain all original edges, and put no edges below \(*\). The relation
\(R^*\) is r.e. and well founded. Code its countable owner set
effectively by \(\mathbb N\).

Applying the theorem to \(R^*\) produces a distinguished agent \(*\)
that knows everyone's truthfulness and mechanicalness, including its
own, and whose particular unary codes are known by every original agent.
Thus such an additional common agent exists for every hierarchy in the
theorem.

A fixed, \(R\)-independent theory for such an agent can be obtained from
the source construction. Apply it to \(R^*\), using a second temporary
auxiliary owner \(\dagger\). The desired \(*\) is now an original owner
relative to \(\dagger\), so \(U_\dagger\) includes mechanicalness for
\(*\).

Choose \(U_\dagger\) and its compiler \(d_\dagger\) once, independently
of \(R\), as in Section 3.2. The least original agent has no
subordinate-code generators. Its enlarged theory is therefore the fixed
theory \[
V_*=
\operatorname{cl}_*\left(
 B_*+
 \left\{\forall x\bigl(
 K_\dagger\varphi\leftrightarrow
 W(\overline{d_\dagger(\varphi)},x)\bigr):
                       \operatorname{FV}(\varphi)=\{x\}\right\}
\right).
\tag{P5}
\] It inherits \(U_\dagger\), including mechanicalness for \(*\). All
other original agents know its particular codes along their added edges.
Removing \(\dagger\) by Section 17.5 gives a fixed distinguished theory
\(V'\) in the desired language.

In fact, \[
\operatorname{Th}(U^+)\subseteq\operatorname{Th}(V').
\tag{P5a}
\] The reduct \(V'\) contains every \(B_*\)-generator, inherits
mechanicalness for every remaining owner from \(U_\dagger\), and is
closed under own sentence necessity. It therefore contains every finite
\(*\)-prefix of every generator of \(U^+\), and consequence closure
gives (P5a).

The constructed standard knowing-machine model interprets \(K_*\) by
\(V'\)-theoremhood and makes every \(V'\)-theorem true. It consequently
satisfies every axiom of \(U^+\). This proves ordinary consistency of
\(U^+\) and gives a standard knowing-machine model of its axiom set.

The distinguished theory \(V'\) need not equal \(U^+\). Replacing its
raw \(*\)-predicate by exact \(U^+\)-theoremhood would require a
separate soundness and genericity argument for that predicate. Formulas
can contain negative occurrences of knowledge, so changing the theory
changes their interpretation; subtheory inclusion alone does not justify
the replacement.

\subsection{18.5. Quantifier order in the finite-room
argument}\label{quantifier-order-in-the-finite-room-argument}

The selected-room statement proved in Sections 12--13 has the following
quantifier order: \[
\forall\theta\ \exists\mathscr R_\theta\;
\forall\psi\ \forall s\ \forall\pi\;
\exists\mathscr D_{\theta,\psi,s,\pi}.
\tag{P6}
\] Here \(\mathscr R_\theta\) is one finite target room and its label
table; \(\psi\) ranges over the target and the formulas in its finitely
many required query-shape classes, including all term variants; \(s\) is
a numerical assignment; \(\pi\) is a finite source certificate for the
required instance; and \(\mathscr D\) is the additional finite proof
room and reflection data. The target table is fixed before every such
variant and instance. The additional proof room may depend on the
certificate.

This is the recovery statement used in the soundness argument: one
target room handles every required query-shape class, while additional
proof rooms may vary with numerical assignments and certificates. No
bound on all certificate depths is needed. The temporary predicates
reconstructed in Section 14 have r.e. numerical columns; a joint
computable selector for their column indices is unnecessary. The
original source family itself remains uniformly r.e. by Section 3.

\subsection{18.6. The boundary for known particular
self-codes}\label{the-boundary-for-known-particular-self-codes}

The theorem gives each agent knowledge of everyone's truthfulness and
mechanicalness, including its own. A stronger particular-code
requirement would give every agent a known particular self-code for
every unary knowledge column while retaining known self-truthfulness.
Reinhardt's negative theorem excludes that full combination under the
relevant Epistemic Arithmetic assumptions {[}A, §1.1, pp.~1--2, 5, and
7{]}. In Alexander's classification, these excluded strengthenings
include Questions 3-2-3 and 3-3-3.

This concerns the entire self-code schema. An agent may know a
particular code for an individual trivial column without contradiction.

\section{References}\label{references}

\textbf{{[}A{]}} Samuel A. Alexander, \emph{The Theory of Several
Knowing Machines}, doctoral dissertation, The Ohio State University,
2013. Relevant locations: §1.1 and Definition 1.1.1 for the
classification and negative boundary; §1.2, Definitions 1.2.3, 1.2.9,
1.2.15, 1.2.17, and 1.2.19, Theorem 1.2.6, and Example 1.2.13 for the
logic and machine semantics; §2.6, p.~75, for inheritance along
transitive closure.
\href{https://etd.ohiolink.edu/acprod/odb_etd/etd/r/1501/10?clear=10&p10_accession_num=osu1366362999}{Official
dissertation record}.

\textbf{{[}W17{]}} Gunnar Wilken, \emph{Tracking chains revisited}, 2017
version. §5.1 and Corollary 5.8 give the elementary-recursive
characterization used for the core.
\href{https://arxiv.org/html/1611.04348v3}{Primary text}.

\textbf{{[}W18{]}} Gunnar Wilken, \emph{Pure patterns of order 2}, 2018
version. Corollaries 4.8 and 4.10 give the finite-pattern realization
and core results used in the finite-tallness argument.
\href{https://arxiv.org/html/1608.08421v5}{Primary text}.

\textbf{{[}W21{]}} Gunnar Wilken, \emph{Pure \(\Sigma_2\)-Elementarity
beyond the Core}, 2021 version. Printed pp.~5--6 include the pure
relations, Lemma 1.7, the core identification, and Theorem 1.8. Theorem
4.2 and Corollary 4.6 supply the corrected semantic characterization.
\href{https://arxiv.org/abs/1710.01870}{Preprint record}.

\end{document}
